\documentclass[11pt,letterpaper]{article}
\usepackage[T1]{fontenc}
\usepackage{lmodern}
\usepackage[margin=.82in,headheight=14pt]{geometry}
\usepackage{amsmath,amssymb,amsthm,mathtools}
\usepackage{booktabs,longtable,array,enumitem,microtype}
\usepackage{xcolor,xurl,hyperref,fancyhdr}
\definecolor{inkblue}{HTML}{173C60}
\hypersetup{colorlinks=true,linkcolor=inkblue,urlcolor=inkblue,citecolor=inkblue,
 pdftitle={Lucas corrections and elimination of arithmetic witnesses},
 pdfauthor={Prepared for Vladimir Reshetnikov with AI assistance},
 pdfsubject={Explicit Diophantine equations for syntactic Con(ZFC)}}
\newcommand{\Z}{\mathbb Z}
\newcommand{\N}{\mathbb N}
\newcommand{\Con}{\operatorname{Con}}
\newcommand{\fl}[1]{\left\lfloor#1\right\rfloor}
\newcommand{\rg}[2]{\{#1,\ldots,#2\}}
\newtheorem{theorem}{Theorem}[section]
\newtheorem{lemma}[theorem]{Lemma}
\newtheorem{proposition}[theorem]{Proposition}
\newtheorem{corollary}[theorem]{Corollary}
\theoremstyle{definition}
\newtheorem{definition}[theorem]{Definition}
\theoremstyle{remark}
\newtheorem{remark}[theorem]{Remark}
\allowdisplaybreaks[3]
\setlength{\emergencystretch}{2em}
\setlength{\parskip}{3pt}
\setlist{nosep,leftmargin=1.5em}
\pagestyle{fancy}
\fancyhf{}
\fancyhead[L]{\small Lucas corrections and witness elimination}
\fancyhead[R]{\small Con(ZFC)}
\fancyfoot[C]{\thepage}
\renewcommand{\headrulewidth}{.3pt}
\pdfinfoomitdate=1
\pdftrailerid{}

\title{\color{inkblue}\bfseries Lucas corrections and elimination\\
of arithmetic witnesses\vspace{5pt}\\
\large Explicit Diophantine forms for $\Con(\mathrm{ZFC})$}
\author{Prepared for Vladimir Reshetnikov and the ProveIt repository\\
\small Research, derivations, code and independent scoped reviews with AI assistance}
\date{25 October 2026\\\small Comparison snapshot:
\texttt{3ecd6a6cf53d47519228328567f45fb7ee9781d9}}

\begin{document}
\maketitle

\begin{abstract}
We give two new arithmetic transformations and a family of explicit polynomial
projections for the Diophantine consistency problem.  A second-order Lucas
correction makes the previously obstructed scale $Y=qs$ usable in the compact
native/Pell construction.  It yields one equation in ten strictly positive
integer variables, of exact degree $304$, evaluated in $87$ charged binary
arithmetic operations.  A separate sign-recovery argument combines two front
nonnegativity certificates in the direct C222 construction.  Its principal
quartic has $248$ unrestricted integer variables, $5{,}321$ collected monomials,
and coefficient height $128{,}446$, improving its matched $252$-variable parent
without increasing degree or height.  Acyclic monic elimination then gives,
among other points, $153$ unrestricted integers at degree $8$, $125$ at degree
$16$, and $102$ at degree $256$.  The complete octic has $188{,}579$ monomials
and height $750{,}273{,}658{,}864$.

For each delivered equation, the new ordinary arithmetic proof, composed with
the identified inherited refutation calibration, proves in ZFC that absence
of solutions in its stated domain is equivalent to syntactic
$\Con(\mathrm{ZFC})$.  The archive supplies literal equations, all fixed
coefficient bytes and their reconstruction, restoration certificates,
independent exact checkers, and complete printed specifications of the
principal compact polynomial and the projected integer family.  Global
minimality and a new end-to-end proof-assistant formalization are not claimed.
\end{abstract}

\tableofcontents
\clearpage

\section{The result, its domain, and its comparison boundary}
\label{sec:overview}

The objective is a fixed polynomial $P$ and a specified domain $D$ such that
\begin{equation}
 \mathrm{ZFC}\vdash
 \bigl(\neg\exists\mathbf x\in D\;P(\mathbf x)=0\bigr)
 \longleftrightarrow\Con(\mathrm{ZFC}).
 \label{eq:target}
\end{equation}
Here $\Con(\mathrm{ZFC})$ denotes the arithmetical assertion that there is no
finite syntactic refutation in the repository's fixed recursive presentation
of ZFC.  Equivalence of reasonable effective proof presentations can be
proved by the corresponding syntactic translations.  The present paper uses
the inherited presentation consistently and does not select coefficients on
the basis of whether it is consistent.

The proof strategy establishes the positive existential statement first:
an allowed zero exists if and only if a finite refutation exists.  Negating
that equivalence gives \eqref{eq:target}.  No assumption of consistency is
used to prove the equivalence, and this paper does not establish either its
unsolvability side or its consistency side.

\subsection{Three useful outcomes rather than an unspecified optimum}

Different measures of simplicity conflict.  An expression with a small
number of unknowns can hide large fixed coefficients; expanding an expression
can change its apparent length without changing the polynomial; eliminating
an arithmetic witness can increase degree.  We therefore give a primary
compact result and a separate family with small explicit circuit numerals.

\begin{center}
\begin{tabular}{@{}lrrp{6.0cm}@{}}
\toprule
Construction & Unknowns & Degree & Main feature\\
\midrule
Compact Lucas correction & $10$ positive & $304$ & $87$ binary operations;
the complete coherent fixed native tuple is supplied.\\[3pt]
Principal direct quartic & $248$ integer & $4$ & $5{,}321$ terms;
height $128{,}446$; direct C222 proof calculus.\\[3pt]
Direct octic projection & $153$ integer & $8$ & Entire expansion supplied:
$188{,}579$ terms, height $750{,}273{,}658{,}864$.\\[3pt]
Moderate-degree projection & $125$ integer & $16$ & Literal circuit and
integral restoration; coefficient height below $2^{80}$.\\[3pt]
Smaller-dimension projection & $102$ integer & $256$ & Literal circuit and
integral restoration; coefficient height below $2^{691}$.\\
\bottomrule
\end{tabular}
\end{center}

``Positive'' means strictly positive, and ``integer'' means unrestricted
integer.  These domains are part of the theorem.  The straight-line
abbreviations in every displayed equation are definitions, not extra
existential variables.  The integer projections have only addition,
subtraction, multiplication and fixed integer constants.  Acyclic
restoration is a finite polynomial construction, not an invocation of a
universal Diophantine polynomial with unspecified parameters.

For the compact equation, each live binary addition, subtraction and
multiplication is charged once; a square and multiplication by a fixed
constant each count as one multiplication.  This is arithmetic circuit cost,
not bit complexity.  The largest of its four stored fixed coefficient ports
has $592{,}884$ bits.  In the direct projection circuits the largest literal
constant is $256$; this does \emph{not} bound the coefficients after expansion.
The separate height columns refer to those collected coefficients.

\subsection{What improves on the inspected sources}

All comparisons are to the pinned repository snapshot above, including the
eight incoming archives recorded in \texttt{SOURCE\_LEDGER.json}.
The latest incoming material matters: comparing only against an older
maintained page would overstate the improvement.  The inspected compact
sources already contain $75/626$, $76/464$ and $77/460$ operation/degree
points, as well as the lower-degree points below
\cite{centered,review,graphdigits,forms}.

\begin{center}
\begin{tabular}{@{}lrrl@{}}
\toprule
Ten-positive construction & Operations & Degree & Status\\
\midrule
Centered residues & 75 & 626 & Prior comparison\\
Centered residues & 76 & 464 & Prior comparison\\
Centered residues & 77 & 460 & Prior comparison\\
Unshifted doubled product square & 84 & 328 & Prior comparison\\
Graph-digit center & 87 & 323 & Prior comparison\\
Sharp graph-digit center & 88 & 318 & Prior comparison\\
Second-order Lucas correction & \textbf{87} & \textbf{304} & This paper\\
Uncentered transport & 89 & 288 & Prior comparison\\
Lower-degree native form & 121 & 246 & Prior comparison\\
\bottomrule
\end{tabular}
\end{center}

Thus the new $87/304$ point improves both displayed graph-digit points in
the cost--degree comparison, with the same ten-positive domain.  It does not
dominate every earlier compact point.  The meaningful mathematical change
is that a scale explicitly rejected in the inspected denominator-engineering
notes becomes possible by correcting the index congruence
\cite{denominator}.  No literature-wide priority claim is inferred from that
repository comparison.

The direct quartic's matched parent has $252$ unrestricted integer variables,
$182$ residuals, $5{,}340$ terms and height $128{,}446$
\cite{common}.  The new $248$-variable principal improves variable and term
counts while retaining the other metrics.  Two further quartics use $246$
and $245$ variables with different height/sparsity tradeoffs.  These counts
should not be presented as global records for unrestricted integer
Diophantine consistency: the repository already contains much smaller
native-tuple forms \cite{forms}.  The distinction here is the direct-calculus
route, its explicit modest coefficients, and the complete algebraic maps.

\subsection{The exact inheritance boundary}

The compact contribution changes the arithmetic front of the existing
native/Pell route.  It retains the complete directly indexed BFS943 tuple
and its original native/refutation calibration.  Its source data ultimately
descend from the existing machine/tile encoding; this paper does not claim
that compact route is independent of K215.  In contrast, the direct quartic
and its higher-degree projections encode the C222 numerical proof calculus.
Their arithmetic does not simulate K215 and does not use Jones's universal
polynomial or Sun's reduction.

For the direct route the inherited final bridges are
\begin{align*}
 &\texttt{KEndpoint.c222\_iff\_zf\_refutation},\\
 &\texttt{ConstructibleChoice.zf\_refutation\_iff\_zfc\_refutation}.
\end{align*}
The first theorem is about the machine-independent calculus even though its
source file also contains machine endpoint theorems.  The second is the
constructibility-based equiconsistency bridge from ZF to ZFC.  The delivered
source closure retains these interface files and identifies their upstream
dependencies.  It is not a new standalone build of all transitive Lean
imports.

The new contributions are ordinary unbounded mathematical proofs plus finite
exact algebraic certificates.  Finite checking authenticates the literal
polynomials, restorations, degree witnesses, coefficient bytes and operation
counts.  It does not replace the Pell, CRT, native-compiler or syntactic
calibration theorems.  A claim of a new end-to-end Lean theorem, or a new
internal-ZFC proof term, would therefore be stronger than this deliverable.
All ``independent reviews'' supplied here are separate checks within an
AI-assisted workflow; they are not outside human peer review.

% BEGIN compact.tex
\section{A second-order Lucas correction: ten positive variables, degree 304}
\label{sec:lucas-q1}

The principal result of this section is a fixed integer polynomial in ten
strictly positive integer variables, with exact degree $304$ and an evaluator
using $87$ binary arithmetic operations.  The arithmetic change repairs a
specific failed direction in the latest inspected submissions: with the
smaller scale $Y=qs$, the elementary first-Pell congruence visibly supplies
only $q^3$, whereas the unknown rank is bounded by $q^4$.  Retaining the next
term of the Lucas polynomial supplies enough precision to recover the index.
The full denominator and the omitted native divisibility are then recovered
in that order.

This is a new step within the inherited native/Pell encoding.  The proof
below composes it with the explicitly retained ordinary calibration of that
encoding to finite syntactic ZFC refutability.  The complete formula, fixed
coefficient bytes and their finite reconstruction are supplied.  No new
end-to-end proof-assistant object or proof of consistency is claimed.

\subsection{The complete polynomial}

The only variables are
\[
 J,F,\alpha,i,T,w,s,n,\eta,\zeta\in\mathbb Z_{>0}.
\]
The four fixed integers $\beta,K_0,M_C,M_F$ are the entire directly indexed
BFS943 tuple of the retained parent.  They satisfy
\[
 \beta\ge31\text{ odd},\qquad K_0>0,\qquad
 0<M_C<\beta<M_F<2\beta.
\]
The largest stored port has $592{,}884$ bits.  These are fixed coefficients,
not additional existential variables.  Their complete finite recipe is
retained separately from the mathematical abbreviations below.

Define the following polynomial expressions:
\begin{gather*}
 h=\beta J,\quad q=h+1,\quad C=q-F-\alpha,\quad X=wq,\quad Y=qs,\\
 k=\eta+\zeta,\quad c=kY+\eta,\quad
 A=Y(X+1)+1,\quad \Delta=A^2-1,\quad H=4A-5,\\
 v=XY^2,\qquad E=1+v(v+1)k^2,\\
 R=[h(q-F)+\alpha](q^2-1)+(M_C+qM_F)J,\qquad
 B=(K_0+w)C-F,\\
 U=ic^2,\quad a=\Delta U^2E,\quad F_0=1+a,\quad G=EF_0,\\
 F_*=1+4aF_0,\quad Q=\Delta(F_*-1),\quad
 b=cT-R,\quad V=F_*b-c,\quad K=X(3c-2X)+2.
\end{gather*}
The new corrected numerator is
\begin{equation}
 e_6=6(k-R-1)-v(R+1)(R-1)(R+3).
 \label{eq:lucas-e6}
\end{equation}
For the principal equation put $L=bV$, and define
\begin{align*}
 D_6&=6q^6,\qquad M=hD_6H,\\
 Z&=hD_6\bigl[H(n-LB)+K\bigr]
       -H\bigl[D_6B+h(e_6-6q^4s)\bigr],\\
 R_G&=M^2G.
\end{align*}
The requested single equation is
\begin{equation}
 \boxed{\;P=Q(Z^2-R_GV^2)+(R_G-Z^2)=0.\;}
 \label{eq:lucas-principal}
\end{equation}
Every named expression is a polynomial abbreviation.  No division, factorial,
binomial coefficient, variable exponent or implicit additional witness occurs
in the equation.  Its literal evaluator is \texttt{equation\_bV.json}.

Replacing $L$ by $V^2$ gives a control equation with degree $410$ in $86$
operations.  It is useful for separating the index correction from the sharper
transport estimate, but is dominated by earlier cost--degree points and is
not claimed as a frontier improvement.

\begin{theorem}[Second-order compact consistency equation]
\label{thm:lucas-q1}
For the complete fixed coefficient tuple just specified, the ordinary
arithmetic proof below, composed with the retained native/refutation
calibration, establishes
\[
 \mathrm{ZFC}\vdash
 \left(\neg\exists\mathbf x\in\mathbb Z_{>0}^{10}\;
 P(\mathbf x)=0\right)\longleftrightarrow\operatorname{Con}(\mathrm{ZFC}).
\]
The principal polynomial has exact degree $304$.  Its supplied evaluator has
$53$ multiplications, $18$ additions and $16$ subtractions, for $87$ binary
operations in total, including squares and multiplication by fixed numerals.
\end{theorem}

For the proof only, write
\[
 t=\frac Bh,\quad W=1+hL,\quad \kappa=\frac KH,\quad
 \sigma=\frac{s}{q^2},\quad \bar e=\frac{e_6}{6},\quad
 \rho=\frac{\bar e}{q^6},\quad z=\frac ZM.
\]
These rational quantities need not initially be integral.  Direct polynomial
identities give
\begin{equation}
 P=(Q-1)Z^2-M^2G(QV^2-1),\qquad
 z=n+\kappa+\sigma-Wt-\rho.
 \label{eq:lucas-carrier}
\end{equation}
In particular the appearance of $1/6$ in the proof does not introduce a
rational-coefficient gate: the circuit computes the integer polynomial $e_6$.

\subsection{Roots and early ranks without assuming source signs}

Positivity gives $q\ge32$, $X,Y\ge q$, $k\ge2$, $c>kY$, $A\ge3$,
$U\ge c^2$, $F_*>2c$, $Q>1$ and $M>0$.  For every integer $b$,
$V=F_*b-c$ is nonzero and $|V|>1$.  The exact identities
\[
 F_*=(1+2a)^2,\qquad Q=4\Delta^2U^2G
\]
show that $G$ is coprime to both $Q-1$ and $QV^2-1$.  At a zero of $P$,
the first identity in \eqref{eq:lucas-carrier} yields
$(Q-1)z^2=G(QV^2-1)$.  Every prime valuation of $G$ is even, so
$G=g^2$ with $g>0$ integral.  Since $\gcd(E,F_0)=1$, its positive factors
are separately squares:
\[
 E=\tau^2,\qquad F_0=f_0^2,\qquad g=\tau f_0.
\]
The positive rational $y=|z|/g$ satisfies
\begin{equation}
 QV^2-(Q-1)y^2=1.
 \label{eq:lucas-adjacent}
\end{equation}

The adjacent-Pell descent forces $y$ integral.  To recall its short proof,
write $y=r/m$ in lowest terms.  Then $m^2\mid Q-1$.  If $u=|V|>1$, then
$u<y<Qu/(Q-1)$, and the norm-preserving transformation
\[
 u'=(2Q-1)u-2(Q-1)y,\qquad
 y'=(2Q-1)y-2Qu
\]
has a nonzero integer first coordinate with $|u'|<u$ and preserves the
reduced denominator $m$ of its second coordinate.  Reflection and descent
reach $u=y=1$, forcing $m=1$.  Thus $z=\varepsilon gy$ is integral,
where $\varepsilon\in\{-1,1\}$.

Set $f=1+2a$ and $I=2Ug$.  Pell doubling gives
\begin{equation}
 f^2-\Delta I^2=1,\qquad Q=(\Delta I)^2,\qquad c^2\mid I,
 \qquad V=f^2(cT-R)-c.
 \label{eq:lucas-strong}
\end{equation}
The retained signed all-index folding theorem applies to
\eqref{eq:lucas-adjacent}--\eqref{eq:lucas-strong}, $c\le I$ and
$V\equiv-c\pmod f$.  It requires no sign for $V$ or $R$ and gives
$c=\psi_A(p)$ for a positive odd $p$.  Here $\psi_A(0)=0$,
$\psi_A(1)=1$ and $\psi_A(j+2)=2A\psi_A(j+1)-\psi_A(j)$.
The first norm gives
\[
 k=2\psi_{2v+1}(n_0),\qquad n_0\ge1.
\]
If $n_0=1$, then $k=2$, $\eta=\zeta=1$, and $c=2Y+1$ lies strictly
between the first two main ordinates $1$ and $2A$.  This is impossible.
Consequently
\begin{equation}
 n_0\ge2,\quad k\ge8v+4,\quad
 c>kY>8XY^3\ge8q^4,\quad c>X,\quad K>0.
 \label{eq:lucas-early-growth}
\end{equation}
No sign of $B,C,R,b,V$ or $\bar e$ has been assumed.

\subsection{Absolute transport at the smaller scale}

If $b=0$, equation \eqref{eq:lucas-adjacent} would give
$y^2-c^2=(c^2-1)/(Q-1)$ strictly between zero and one.  Hence $b\ne0$,
$\operatorname{sgn}V=\operatorname{sgn}b$, and $|V|>|b|$.  With
$D_0=\Delta U^2$, one has
\[
 \frac{a^2}{g^2}=\frac{D_0a}{1+a}>\frac{D_0}{2}>9.
\]
Thus $a>3g$, $|V|>3g$ and $y<3|V|/2$.  In particular
$gy<V^2/2$ for the control.

For $L=bV$, the third main index must also be excluded.  Suppose
$c=4A^2-1$.  If $n_0\ge3$, then $k>32v^2$, while $A<2XY$ gives
\[
 c<16X^2Y^2<32X^2Y^5<kY<c,
\]
a contradiction.  Thus $n_0=2$ and $k=8XY^2+4$.  Put
$u=X+2-2Y$.  Direct expansion gives
\[
 \eta=4Y^2(Xu+1)+8Y(X+1)+3-4Y.
\]
For $u\le-1$ this is negative, using $Y(X-1)>2(X+1)$.  For $u\ge2$,
subtracting $k$ leaves at least $4Y^2+8Y(X+1)-4Y-1>0$.  Both contradict
$0<\eta<k$.  The remaining $u=0,1$ would make $2Y-X$ equal to $2$ or $1$,
contradicting $q\mid X,Y$ and $q\ge32$.  Therefore $p\ge5$ and
$c>c_3:=4\Delta+3$.

Put $S=\Delta I$, so $S^2=Q$.  Pell classification of
\eqref{eq:lucas-adjacent} gives $|V|=\chi_S(j)/S$ for a positive odd $j$.
Index one gives $|V|=1$ and is excluded.  At index three,
$|V|=4Q-3$ would give one of
\[
 F_*(b-4\Delta)=c-c_3,\qquad F_*(b+4\Delta)=c+c_3.
\]
Both right sides lie strictly between zero and $2c<F_*$, impossible.
Hence $j\ge5$, $|V|\ge16Q^2-20Q+5>8Q^2$, and
\[
 |b|>\frac{8Q^2-c}{F_*}>4F_*>4g.
\]
It follows that $L=bV=|b||V|>0$ and $gy<3L/8<L/2$.  For either transport
choice we have established
\begin{equation}
 L>0,\qquad gy<L/2,\qquad |b|/q^6<L/2.
 \label{eq:lucas-transport-bounds}
\end{equation}

\subsection{The corrected numerator and all possible rank signs}

The new correction needs a fresh sign argument: its cubic changes sign at
negative ranks.  Set $e_0=k-R-1$ and
$B_R=(R+1)(R-1)(R+3)$.  Since $b-e_0=cT-k+1>0$,
\[
 \rho<\frac b{q^6}-\frac{vB_R}{6q^6}.
\]
When $B_R\ge0$, this implies $\rho<|b|/q^6$.  If $B\le-1$, then
$Wt<-L$ and \eqref{eq:lucas-carrier}, with positive $\kappa,\sigma$ and
\eqref{eq:lucas-transport-bounds}, gives $n<0$.

The integer exceptions are exactly $R\le-4$ and $R=0$.  Write
$P_{\rm mask}=(M_C+qM_F)J>0$.  The exact outer identity is
\begin{equation}
 R=(hC+q\alpha)(q^2-1)+P_{\rm mask}.
 \label{eq:lucas-signed-outer}
\end{equation}
For $R=-r\le-4$, it implies $C<0$ and
\[
 |C|>\frac r{h(q^2-1)}>\frac r{q^3}.
\]
Thus $B=-(K_0+w)|C|-F\le-1-w|C|$ and
$Wt<-(1+w|C|)L$.  The adverse correction is bounded by
\begin{equation}
 \delta_0=\frac{v(r-1)(r+1)(r-3)}{6q^6}
          <\frac{ws^2r^3}{6q^3},
 \label{eq:lucas-adverse}
\end{equation}
because $v=wq^3s^2$.  Here $b=cT+r>r,c$, so
$V>(F_*-1)b$ and either choice of $L$ satisfies
$L>(F_*-1)r^2$.  Also $\Delta>Y^2\ge s^2$ and
$F_*-1=4a(1+a)>s^2$.  Consequently $w|C|L>\delta_0$.  Substitution gives
\[
 n<L/2-(1+w|C|)L+|b|/q^6+\delta_0<0.
\]

For $R=0$, \eqref{eq:lucas-signed-outer} again gives $C<0$, while
$\delta_0=v/(2q^6)=ws^2/(2q^3)$.  Now $b=cT\ge c$ and
$V\ge(F_*-1)b$, so $L>F_*-1$ and
$w|C|L>w(F_*-1)>\delta_0$.  The same contradiction follows.
This handles every integer rank, including the cases where the new cubic
correction has the unfavorable sign.

Therefore every positive zero has $B\ge0$, which forces $C>0$ because
$F>0$.  The retained mask bound now applies:
\begin{equation}
 3q<R<q^4-2.
 \label{eq:lucas-rank-box}
\end{equation}
Equation \eqref{eq:lucas-early-growth} gives $c>2R$, whence $b,V>0$.
The retained normalized rank and signed auxiliary-index theorems have all
their original hypotheses: $c^2\mid I$, $c>2R,c>2p$,
$V\equiv-c\pmod f$ and $V\equiv-R\pmod c$.  They give $p=R$.
In particular $R$ is positive odd; the unchanged mask parity gives even $q$
and odd $J$.

\subsection{The second-order index lemma}

The new argument is elementary and works at every prime, including the primes
dividing the displayed denominator $6$.

\begin{lemma}[Corrected index recovery]
\label{lem:lucas-corrected-index}
At every positive zero of \eqref{eq:lucas-principal},
\[
 n_0=(R+1)/2,\qquad q^6\mid\bar e,
 \qquad \rho\in\mathbb Z_{\ge0},\qquad q^2\mid s.
\]
The conclusions hold as well for the square-transport control.
\end{lemma}

\begin{proof}
The Lucas polynomial has the finite expansion
\begin{equation}
 \psi_{1+2v}(m)=\sum_{j=0}^{m-1}(4v)^j
                    \binom{m+j}{2j+1}.
 \label{eq:lucas-binomial-expansion}
\end{equation}
For a direct derivation, its recurrence gives the formal generating function
\[
 \sum_{m\ge1}\psi_{1+2v}(m)t^{m-1}
 =\frac1{(1-t)^2-4vt}
 =\sum_{j\ge0}\frac{(4v)^jt^j}{(1-t)^{2j+2}}.
\]
Coefficient comparison yields \eqref{eq:lucas-binomial-expansion}.  In
particular, for an explicitly determined nonnegative integer $Z_m$,
\begin{equation}
 2\psi_{1+2v}(m)
 =2m+8v\binom{m+1}{3}+32v^2Z_m.
 \label{eq:lucas-second-order}
\end{equation}
The formula also covers $m=1,2$, when $Z_m=0$.

Put $r_0=(R+1)/2$, an integer by the preceding rank argument.  Then
$B_R=48\binom{r_0+1}{3}$, so $\bar e=e_6/6$ is an integer and
\begin{equation}
 \bar e=k-2r_0-8v\binom{r_0+1}{3}.
 \label{eq:lucas-ebar}
\end{equation}
The integrality of $z$ and the carrier give
$6q^6\mid hH(e_6-6q^4s)$.  Cancelling the literal common factor $6$,
now that \eqref{eq:lucas-ebar} is integral, gives
$q^6\mid hH(\bar e-q^4s)$.  Because $h=q-1$ and
$H=4Y(X+1)-1\equiv-1\pmod q$, one has $\gcd(q,hH)=1$.  Hence
\begin{equation}
 q^6\mid\bar e-q^4s,\qquad q^4\mid\bar e.
 \label{eq:lucas-carrier-congruence}
\end{equation}
No inverse of $6$ modulo $q$ has been taken.

Set $d=2n_0-R-1=2(n_0-r_0)$.  Reducing
\eqref{eq:lucas-second-order} and \eqref{eq:lucas-ebar} modulo $4v$
gives $\bar e\equiv d\pmod{4v}$.  Since $q^3\mid v$, equation
\eqref{eq:lucas-carrier-congruence} first gives $q^3\mid d$.
The exact Taylor difference is
\begin{align}
 6(\bar e-d)
 &=vd\bigl[(2n_0)^2+(2n_0)(R+1)+(R+1)^2-4\bigr]
       +192v^2Z_{n_0}.
 \label{eq:lucas-taylor-upgrade}
\end{align}
Both terms are divisible by $q^6$.  If a prime $\ell\mid q$ has
$a=v_\ell(q)\ge1$, then $v_\ell(6)\le1$, so
\[
 v_\ell(\bar e-d)\ge6a-1\ge4a.
\]
Thus $q^4\mid\bar e-d$, uniformly at $\ell=2,3$ as well.  Together
with \eqref{eq:lucas-carrier-congruence}, this gives $q^4\mid d$.

The first base $2XY^2+1$ exceeds $A$.  Were $n_0\ge R$, monotonicity
in the positive base and index would give $k/2\ge\psi_A(R)=c$,
contrary to $c>k$.  Therefore $1\le n_0<R$, and
\eqref{eq:lucas-rank-box} gives $|d|<R<q^4$.  Hence $d=0$.

Only now return to the full expansion: equations
\eqref{eq:lucas-second-order}--\eqref{eq:lucas-ebar} give
$\bar e=32v^2Z_{r_0}$.  Because $q^3\mid v$, this proves
$q^6\mid\bar e$ and $\rho\ge0$ integral.  Finally,
\eqref{eq:lucas-carrier-congruence} gives $q^6\mid q^4s$, hence
$q^2\mid s$.
\end{proof}

The two congruence steps are essential.  The first establishes divisibility
of the index error by $q^3$; this makes the correction error small in the
relevant prime valuations and upgrades it to the $q^4$ modulus needed for
the ordinary size bound.  The full $q^6$ denominator is recovered after the
index has been identified.

\subsection{Projection and restoration of the native source}

Lemma~\ref{lem:lucas-corrected-index} makes $\sigma=s/q^2$ integral and
restores $q^3\mid Y=qs$ before projection.  The carrier now gives
$\kappa-WB/h\in\mathbb Z$, so $H\mid hK$.
Write $R=2r+1$.  By \eqref{eq:lucas-rank-box}, $r\ge48$.
The retained ratio theorem, uniformly valid at scale one, gives
\[
 Y>X^r/3,\qquad X^{r+1}<3Y(X+1)<H.
\]
The universal Pell projector and the retained entire-product estimate give
\[
 H\mid h(X-2^R)(2^RX+1),\qquad
 |h(X-2^R)(2^RX+1)|<H.
\]
The product is zero, so $X=2^R$ and $\kappa$ is integral.  No possibly
nonunit factor of the product has been cancelled modulo $H$.  The remaining
carrier gives $h\mid WB$, and $W\equiv1\pmod h$ gives $h\mid B$.
Thus the native quotient $t=B/h$ is a nonnegative integer.

All inherited fixed-minus parity and exact half-binomial rounding hypotheses
have now been restored.  The fresh reconstruction uses the original source
base $Y(X+1)+2$, which is distinct from the auxiliary base $A$ used here.
It produces an actual zero of the calibrated native source.  The whole fixed
coefficient tuple's retained native/refutation theorem therefore yields a
finite syntactic ZFC refutation from every positive zero of
\eqref{eq:lucas-principal}.

\subsection{Construction of positive witnesses}

Conversely, a finite refutation supplies the parent's canonical outer data:
$q^3\mid Y$, $X=2^R$, $R\equiv3\pmod4$, $R>3q$ and integral
$t=B/h\ge0$, with the exact half-binomial interval.  Define the new positive
input $s=Y/q$.  Then $\sigma=s/q^2=Y/q^3$ is an integer and at most $Y$.
Using the scale-one bases, put
\begin{gather*}
 c=\psi_A(R),\qquad
 k=2\psi_{2XY^2+1}\bigl((R+1)/2\bigr),\\
 \tau=\chi_{2XY^2+1}\bigl((R+1)/2\bigr),\qquad
 \eta=c-kY,\quad\zeta=k-\eta.
\end{gather*}
The inherited uniform ratio estimate gives $\eta,\zeta>0$.  The full
Lucas expansion gives $\rho\ge0$ integral; projection gives positive
integral $\kappa<Xc$.

A positive multiple of a finite Pell matrix period modulo $c^2\tau$
supplies one seed pair $(f_0,I_0)$ with $c^2\tau\mid I_0$.
Set $i=I_0/(c^2\tau)>0$ and double this pair exactly as in the parent.
It yields \eqref{eq:lucas-strong}, $G=(\tau f_0)^2$ and the canonical
positive auxiliary pair $V,y$, with the positive integral unshifted input
\[
 T=\frac{V+c+Rf^2}{cf^2}.
\]
No private strong pair is reselected after its auxiliary congruences are
specified.  Moreover
\[
 f_0>(A-1)I_0>2Xc,\qquad \kappa<Xc,\qquad \sigma\le Y<Xc,
\]
so $g=\tau f_0>\kappa+\sigma$.  Now set
\[
 n=gy+(1+hL)t+\rho-\kappa-\sigma.
\]
This is a strictly positive integer, and the resulting carrier is $z=gy$.
The polynomial therefore vanishes.  Every one of the ten independent
inputs has been constructed positively.  This completes both existential
directions in Theorem~\ref{thm:lucas-q1}.

\subsection{Exact degree and validation}

The relevant global degrees are
\[
 \begin{array}{c|rrrrrrrr}
 \text{expression}&Y&c&A&\Delta&v&E&U&a\\\hline
 \deg&2&3&4&8&6&14&7&36
 \end{array}
 \qquad
 \begin{array}{c|rrrrrrr}
 \text{expression}&F_*&Q&G&b&V&M&Z\ (L=bV)\\\hline
 \deg&72&80&50&4&76&11&93
 \end{array}
\]
The correction $v(R+1)(R-1)(R+3)$ has degree $18$, below the leading
transport part.  For $L=bV$, the term $-QM^2GV^2$ uniquely dominates the
final expression and has degree $304$.  For $L=V^2$, $\deg Z=165$ and
$QZ^2$ uniquely supplies degree $410$.

At the smaller scale, $cT$ and $R$ both have degree four; the leading part
of $b=cT-R$ must retain both.  Put
\[
 q_0=\beta J,\qquad k_0=\eta+\zeta,\qquad C_0=q_0-F-\alpha,
 \qquad B_0=k_0sT-q_0^2(q_0-F).
\]
The entire highest homogeneous polynomials are
\begin{align*}
 [P_{bV}]_{304}
   &=-9\cdot2^{12}q_0^{128}w^{34}s^{78}i^{14}k_0^{44}B_0^2,\\
 [P_{V^2}]_{410}
   &=9\cdot2^{16}q_0^{166}w^{46}s^{104}i^{20}k_0^{60}B_0^4C_0^2.
\end{align*}
For example, their respective coefficients at
\[
 J^{128}w^{34}s^{80}i^{14}\eta^{46}T^2,\qquad
 J^{166}w^{46}s^{108}i^{20}\eta^{64}T^4F^2
\]
are $-9\cdot2^{12}\beta^{128}$ and
$9\cdot2^{16}\beta^{166}$, both nonzero at the authentic fixed tuple.
These give global degree certificates, rather than sampled line degrees.

The independent exact checker binds all actual evaluator rows through $37$
free-polynomial cuts per circuit.  It recounts the live gates, checks every
fixed coefficient byte, and propagates the complete highest homogeneous
forms: $274$ terms for the
principal equation and $3{,}750$ for the control.  Actual operator corruption
checks reject $175$ modified principal artifacts and $173$ control artifacts.
Normal and optimized Python runs agree.  A separate retained reconstruction
script regenerates all four fixed numerals from their finite recipe and
checks them against the stored bytes.  The Taylor identity
\eqref{eq:lucas-taylor-upgrade} is checked exactly; separate finite Lucas
recurrence controls are explicitly diagnostic.

Those finite checks certify the equations and algebraic certificates.  They
do not replace the unbounded arithmetic proof above or the named inherited
calibration, and do not decide the existence of a ZFC refutation.  The new
work has also received a separate ordinary mathematical review within this
AI-assisted workflow; that is not an outside human referee's endorsement.
% END compact.tex

% BEGIN direct.tex
\section{A direct quartic from one certificate for two front signs}
\label{sec:direct-front}

The second contribution is a local arithmetic simplification of the direct
C222 construction.  Its principal equation has 248 unrestricted integer
variables, degree four, 5,321 collected monomials, and coefficient height
128,446.  It is the sum of the squares of the 182 explicit quadratic
residuals in
\texttt{equations/integer248/equation.json}.  Here and below, coefficient
height is measured after collecting the entire polynomial.

\begin{center}
\begin{tabular}{lrrrrr}
Profile & Variables & Residuals & Degree & Terms & Height\\ \hline
Common Remainders parent & 252 & 182 & 4 & 5,340 & 128,446\\
Principal & 248 & 182 & 4 & 5,321 & 128,446\\
Balanced projection & 246 & 180 & 4 & 5,048 & 195,955\\
Dimension projection & 245 & 179 & 4 & 21,845 & 314,701
\end{tabular}
\end{center}

These are three achieved projections of one new raw source, not an
assertion that every better projection has been excluded.  The principal
profile strictly improves the dimension and term count of its matched
parent while preserving degree, residual count, and coefficient height.
The new mathematical step is a proof that one product certificate can
replace two front sign certificates, even though the product alone has
an unwanted negative branch.  The exact factorial digit already present
in the source excludes that branch.

\subsection{Literal source change and inherited interface}

Put
\[
 G(a,b,c)=a^2+b^2+c(c+1).
\]
This is nonnegative on $\Z^3$ and represents every nonnegative integer.
Indeed, the three-square theorem represents $4n+1$ as three squares; its
residues modulo four put exactly one odd square in that representation.
Writing it as $(2c+1)^2$ and the others as $(2a)^2,(2b)^2$ gives $n=G(a,b,c)$.

Use the following semantic names for the unchanged source coordinates:
\[
\begin{split}
 r&=\texttt{fact0.remainder},\quad
 s=\texttt{fact0.slack},\quad F=\texttt{fact0.output},\\
 B&=\texttt{choose0.output},\quad
 V=\texttt{pow0.output},\quad q=\texttt{pow2.output}.
\end{split}
\]
The old front separately imposed $r\geq0$, $s\geq0$, and
$o\geq0$, where $q=B+o+1$.  Replace its two first ternary guards by
\begin{equation}
 (r+1)(s+1)=1+G(a,b,c).
 \label{eq:front-fused}
\end{equation}
This is the quadratic row $rs+r+s-a^2-b^2-c^2-c=0$; it requires no
multiplication gate.  Compose the previously supplied odd-square radix
restriction by deleting $q=B+o+1$ and its private guard, and inserting
\begin{equation}
 q=(2z+1)^2.
 \label{eq:front-odd}
\end{equation}
The two remaining guards---the strict squared Pell cap and the common
normalized remainder norm---are unchanged.  All other source rows,
all eighteen power requests, and all endpoint and Horn-flow relations
are retained.  The old coordinate $o$ occurs only in its deleted bound
row and can therefore be restored freely at the end.

The starting point is the direct component of
\path{ConZFC_Common_Remainders_Radix47_2026-10-25.zip}, at repository
snapshot \path{3ecd6a6cf53d47519228328567f45fb7ee9781d9}.
Its \path{NORMALIZED_REMAINDERS_PROOF.md}, together with the retained
\path{ORIENTATION_FIRST_PROOF.md} and
\path{STRICT_BUDGET_PROOF.md}, supplies the following interface
\emph{before} the removed front signs are used.  An integer $W\geq258$
and a signed Pell center $A$ satisfy
\begin{equation}
 |A|>W^W,\qquad |n|,|k|,|m|<W
 \label{eq:front-frame}
\end{equation}
for every requested base, index, and output.  The corresponding principal
Pell coordinates are canonical at the signed nominal index $k$.
The frame uses squares of the proposed atoms, and the squared semantic
anchor with bounded pooled errors orients both Pell coordinates.
Consequently those conclusions do not presuppose the signs of $F,B,q$
or of the independent CRT representatives.  We retain that inherited
orientation theorem and extend its subsequent power argument to signed
bases explicitly below.

\subsection{The signed-base extension}

\begin{lemma}[Signed bases]
\label{lem:front-signed}
Under the inherited interface \eqref{eq:front-frame}, a power request
with $2\leq |n|<W$ forces $k\geq0$ and $m=n^k$.  The assertion holds
for either sign of $n$ and either sign of $A$.
\end{lemma}
\begin{proof}
Set $T=2An-n^2-1$.  The actual output row is
\[
 m\equiv X_k(A)+(n-A)Y_k(A)\pmod T.
\]
The canonical signed pairs have
$X_{-j}(A)=X_j(A)$ and $Y_{-j}(A)=-Y_j(A)$.
Furthermore
\[
 |T|\geq2|A||n|-n^2-1
       \geq4|A|-(W-1)^2-1>W^W+W,
\]
where the last inequality follows from $W\geq258$ and $|A|>W^W$.
The absolute value on $A$ is essential: no positive-center premise has
been introduced.  Since $T\equiv-1\pmod n$, $n$ is invertible modulo $T$.

Substituting $n-A$ into the Pell identity gives
$(n-A)^2\equiv A^2-1\pmod T$.  The canonical Pell addition identities
therefore give $m\equiv n^k\pmod T$ when $k\geq0$.  In that case
\[
 |m-n^k|\leq |m|+|n|^k<W+W^W<|T|,
\]
so the congruence is an equality.  If $k=-j<0$, the signed pair is the
inverse of the index-$j$ Pell unit; its norm is one.  Thus
$mn^j\equiv1\pmod T$.  Now $1\leq j\leq W-1$ and $|m|\leq W-1$, so
\[
 |mn^j-1|\leq(W-1)^W+1<|T|.
\]
It follows that $mn^j=1$, impossible in $\Z$ because $|n|^j\geq2$.
Hence $k\geq0$ and the first argument applies.
\end{proof}

For the separate base-one case, a known nonnegative index gives
$m\equiv1\pmod{2A-2}$; the same size argument gives $m=1$.
This case is needed because $F=0$ has not yet been excluded.  The known
index $t\geq16$ in the request at base $F^2+1$ consequently gives
\begin{equation}
 V=(F^2+1)^t\geq1
 \label{eq:front-V}
\end{equation}
for every initially possible integer $F$.

\subsection{Both front branches have the needed size bound}

The untouched ceiling relations are
\begin{equation}
 r+s=B-1,\qquad BF=V+s.
 \label{eq:front-ceiling}
\end{equation}
The right side of \eqref{eq:front-fused} is positive.  Since its two
factors are integers, it leaves precisely the alternatives
\[
 \text{(I) }r,s\geq0,
 \qquad
 \text{(II) }r,s\leq-2.
\]
In case I, $B\geq1$ and $BF=V+s\geq1$ force $F\geq1$.
Since $V=(F^2+1)^t>F^2$, the inequality $V\leq BF$ then gives
$B>F$.  In particular $B\geq2$ and
\[
 V\leq BF<B^2.
\]

In case II put $B=-b$, so $b\geq3$.  If $F=0$, the base-one case of
\eqref{eq:front-V} gives $V=1$, and $BF=V+s$ would give $s=-1$.
If $F\geq1$, combining \eqref{eq:front-ceiling} gives
\[
 r=(1-F)B+V-1>0,
\]
because $(1-F)B\geq0$ and $V>1$.  Both possibilities contradict
case II.  Therefore $F=-f$ for some $f\geq1$.  The same two ceiling
identities now give
\begin{equation}
 bf+2\leq V\leq(f+1)b-1<(f+1)b.
 \label{eq:front-negative-bound}
\end{equation}
For $f\geq1$ and $t\geq16$, $(f^2+1)^t>(f+1)^2$.
Comparing this with \eqref{eq:front-negative-bound} shows $b>f+1$,
and hence $V<b^2$.

Thus, in both branches, defining $f=|F|$ and $b=|B|$ gives
\begin{equation}
 f\geq1,\qquad b\geq2,\qquad V=(f^2+1)^t<b^2.
 \label{eq:front-uniform}
\end{equation}
This deduction has used no factorial fixed point or binomial-digit
interpretation.  That separation is what makes the later sign recovery
noncircular.

\subsection{Recovering the radix before resolving the branch}

Equation \eqref{eq:front-odd} first gives $q\geq1$, so $q+1\geq2$.
Each of the nine actual CRT representatives $a_i$ is an exponent in a
retained request at base $q+1$.  Lemma~\ref{lem:front-signed} therefore
makes each $a_i\geq0$ and authenticates its output.  Set
\[
 H_0=1+\sum_{i=1}^{9}a_i\geq1.
\]
These are actual representatives; several delivered coordinates are
shifted by one and need not themselves be nonnegative.

Apply the lemma next to the retained request $N=B^x$.  Its base has
$|B|\geq2$ by \eqref{eq:front-uniform}, so it forces $x\geq0$ even
while $q=1$ is still an initial possibility.  The unchanged cofinal beta
choice is
\[
 y=V(V-1)>f+2.
\]
In unshifted semantic coordinates, the retained endpoints are
\begin{equation}
 x=F+1+(y+1)u_0,
 \qquad
 x=1+\bigl(1+(t+1)y\bigr)u_t.
 \label{eq:front-endpoints}
\end{equation}
The second equality excludes $x=0$ because its bracket is an integer
greater than one.  If $x=1$, the first would say
$F+(y+1)u_0=0$, impossible since $0<|F|<y+1$.  Consequently $x\geq2$,
and the second equality then forces $u_t\geq1$.  Thus
\begin{equation}
 x\geq2+(t+1)y\geq8.
 \label{eq:front-x-bound}
\end{equation}
The coordinate translations used in the emitted polynomial do not alter
these semantic identities; their explicit shifts are recorded separately.

The request $q=(BH_0)^x$ now has $|BH_0|\geq2$ and is authenticated
by the lemma.  Since its output is positive, $B<0$ would force $x$ even.
Therefore in either branch
\begin{equation}
 N=b^x>0,\qquad q=b^xH_0^x\geq b^x,\qquad N\mid q.
 \label{eq:front-radix-bound}
\end{equation}
In particular $q>1$.  All denominator requests are now genuine positive
powers of $q$.  Write
\[
 M=q^x,\qquad D_0=q^t,\qquad D_1=q^{t+1},\qquad D_y=q^{y+1}.
\]
By \eqref{eq:front-x-bound}, these three denominators divide $M$.
Nothing in this argument has assumed $F=t!$ or the sign of $B$.

\subsection{The exact factorial digit chooses the positive branch}

For $m\geq1$, $t\geq1$ and $0\leq j\leq t$, the elementary estimate
\begin{equation}
 \binom{m+t}{j}\leq(m+1)^t
 \label{eq:front-binomial-bound}
\end{equation}
holds.  One proof is induction on $t$: Pascal's identity bounds each
nonedge entry in the next row by $2(m+1)^t\leq(m+1)^{t+1}$; the new
edge is bounded by
$\prod_{i=1}^{t+1}(1+m/i)\leq(m+1)^{t+1}$.
Apply this with $m=F^2\geq1$.  Every coefficient through position $t$
of $(q+1)^{F^2+t}$ is at most $V<b^2$.  Its true low remainder at
position $t$ is therefore
\[
 r^0=\sum_{j=0}^{t-1}\binom{F^2+t}{j}q^j,
 \qquad
 0\leq r^0<\frac{V D_0}{q-1}<\frac{D_0}{32}.
\]
For the last inequality, \eqref{eq:front-radix-bound} and
\eqref{eq:front-x-bound} give
\[
 \frac{V}{q-1}<\frac{b^2}{q-1}
 <\frac{2}{b^{x-2}}\leq\frac1{32}.
\]
Thus the displayed sum really is the ordinary remainder modulo $D_0$;
the polynomial tail is divisible by $D_0$.

Here is the precise role of the inherited common norm.  Its normalized
digit rows have $D_i v_i=M$ and
$v_iU_i=M Q_i+R_i$.  Since $D_i,M>0$, $v_i$ is a positive integer.
The digit row makes $v_i\mid R_i$, so $r_i=R_i/v_i$ is an integer.
The norm equation
\begin{equation}
 M^2=1+4\sum_{i=1}^{11}R_i^2+G(\alpha,\beta,\gamma)
 \label{eq:front-norm}
\end{equation}
gives $|R_i|<M/2$, hence $|r_i|<D_i/2$.
For the factorial row, $r_i\equiv r^0\pmod{D_0}$.
Together with $0\leq r^0<D_0/32$, this forces $r_i=r^0$.
No sign of $B$ is used in this recovery.

The exact factorial digit row is now
\[
 (q+1)^{F^2+t}=D_0(B+qh)+r^0.
\]
Divide by $D_0$ after subtracting $r^0$ and reduce modulo $q$ to obtain
\begin{equation}
 B\equiv C\pmod q,
 \qquad C=\binom{F^2+t}{t}>0.
 \label{eq:front-final-digit}
\end{equation}
Unlike the old proof, this argument does not yet assume $0\leq B<q$.
Instead, \eqref{eq:front-uniform} and the coefficient estimate give
\[
 |B-C|\leq b+C\leq b+V<b+b^2<q,
\]
the last inequality following from $q\geq b^x$, $b\geq2$ and $x\geq8$.
Equation \eqref{eq:front-final-digit} therefore gives $B=C>0$.
Case II is impossible.  We have recovered $r,s\geq0$, $F\geq1$,
and $q>B$.  Restore the old private slack by $o=q-B-1\geq0$ and fill
all three deleted ternary guards using surjectivity of $G$.

At this point the previous Common Remainders source has an integer
solution.  Its exact factorial fixed-point theorem gives $F=t!$.
Its remaining ten low-digit estimates, prime-power CRT range decoding,
and Horn-flow interpretation therefore apply with their original
premises.  The new argument has restored those premises in the order
needed: size bound, radix, normalized remainder, factorial digit, front
signs, then factoriality.

\subsection{Completeness and the parity restriction}

For the converse, begin with a finite C222 refutation and use the
retained cofinal arithmetic construction.  The odd-square radix
restriction is inherited from the separate Centered Residues proof,
included as \texttt{INHERITED\_ODD\_SQUARE\_PROOF.md}; its compatibility
with the present construction can be seen explicitly.

After factoriality, $F=t!$ is even.  Each $i\leq t$ divides $F$ and
$F^2/i$ is even, so
\[
 B=\binom{F^2+t}{t}=\prod_{i=1}^t\left(1+\frac{F^2}{i}\right)
\]
is odd.  Also $y=V(V-1)$ is even, so every beta CRT modulus is odd.
The CRT progression for $x$ consequently has odd period, allowing $x$
to be chosen even and as large as required.  Now $N=B^x$ and $N^2$
are odd.  The independent amplified CRT representatives can be shifted
by multiples of $N^2$ without changing their required residues.
Adjusting a third representative's parity makes $H_0$ odd while leaving
the two representatives used for the dominant anchor's parity free to
retain opposite parities.  The dominant lift can still be made large
enough for all strict bounds.  Thus $BH_0$ is odd and $x$ is even;
$q=(BH_0)^x$ is the square of an odd integer and
\eqref{eq:front-odd} has an integer witness $z$.

The genuine old front has $r,s\geq0$, so $rs+r+s\geq0$ has a
$G$ representation and satisfies \eqref{eq:front-fused}.  The old
cofinal low-remainder estimates give $0\leq r_i^0<D_i/32$ for all
eleven digits.  Setting $R_i=(M/D_i)r_i^0$ gives $R_i<M/32$ and
\[
 4\sum_{i=1}^{11}R_i^2
 <\frac{44}{1024}M^2<M^2.
\]
Hence the nonnegative integer
$M^2-1-4\sum_iR_i^2$ has a $G$ representation, preserving
\eqref{eq:front-norm}.  Every other retained row is the old row with
its cofinal witness.  This constructs a zero of the new raw source.

The raw modification therefore preserves \emph{existence} of integer
zeros in both directions.  It is not a bijection on all arbitrary old
private-witness tuples.  The later monic projections and integral
translations, in contrast, have explicit reverse-restoration maps.

\subsection{Exact endpoint and the logical calibration chain}

\begin{theorem}
\label{thm:direct248}
Let $P_{248}\in\Z[x_0,\ldots,x_{247}]$ be the sum of the squares of
the 182 delivered residuals in the principal profile.  With the
retained direct-C222 arithmetic and logical calibration,
\[
 \mathrm{ZFC}\vdash
 \left(\exists x\in\Z^{248}\;P_{248}(x)=0\right)
 \longleftrightarrow
 \text{there exists a finite ZFC refutation}.
\]
Consequently ZFC proves that the absence of integer zeros of $P_{248}$
is equivalent to $\operatorname{Con}(\mathrm{ZFC})$.
The same conclusions hold for the two delivered 246- and 245-variable
quartics.
\end{theorem}
\begin{proof}
Over the integers a finite sum of squares vanishes exactly when every
residual vanishes.  The supplied integral translations are bijections;
each monic elimination has an integer polynomial restoration, replayed
in reverse order.  The independently checked Horn factoring identities
preserve the source's solution existence.  The preceding soundness and
completeness arguments identify zeros of the new source with zeros of
the Common Remainders source.  The retained arithmetic proof decodes
those zeros into finite C222 refutations, and constructs arithmetic
witnesses from such refutations.

The last two logical bridges are the identified repository theorems
\path{KEndpoint.c222_iff_zf_refutation} and
\path{ConstructibleChoice.zf_refutation_iff_zfc_refutation}.
Thus the complete chain is: explicit quartic zero, new raw arithmetic
witness, old Common Remainders witness, finite C222 refutation, finite
ZF refutation, and finite ZFC refutation.  Every arrow used here is
either the new arithmetic argument above, a finite exact algebraic
certificate, or the named inherited calibration.  The first-order
arithmetical statement that no finite syntactic ZFC refutation exists
is $\operatorname{Con}(\mathrm{ZFC})$; negating the existence
equivalence gives the stated result.
\end{proof}

The package retains all 88 files in the parent's arithmetic and proof
closure, pinned byte-for-byte.  Its new finite verifier imports the
retained independent algebra checker, not the producing eliminator.  It
reconstructs the complete raw mutation, checks 18 Horn-factor edits,
replays every monic restoration and translation, and expands each full
sum of squares.  The principal profile has 203 monic elimination steps
and 70 translated coordinates.  Independent runs under ordinary Python
and under \texttt{python -O} pass, including five rejected deliberate
mutations; a clean rebuild reproduces all 22 generated JSON files
byte-for-byte.  These are algebraic checks of the emitted artifacts.
The unbounded arithmetic proof remains an ordinary mathematical proof,
with a separate independent review; no new end-to-end Lean elaboration
or proof of consistency is being claimed.
% END direct.tex

% BEGIN elimination.tex
\section{Trading degree for fewer direct proof coordinates}
\label{sec:monic-elimination}

The preceding sign bootstrap gives a direct equation in $248$ unrestricted
integers.  We now reduce its dimension by allowing higher degree.  This
route retains the direct C222 proof calculus: no K215 simulation, Jones
universal equation, or large fixed coefficient tuple is introduced.
The elimination principle is standard.  The contribution here is the
complete family of explicit polynomials and the certified tradeoffs below.

Write the preceding source as
\[
 P_{248}(x)=\sum_{j=0}^{181}R_j(x)^2,
 \qquad x=(x_0,\ldots,x_{247})\in\Z^{248}.
\]
The complete quadratic residual list is retained without change in
\texttt{source/equation248.json}; its source hash is recorded in
\texttt{SOURCE.json}.  The pivot tables in
Appendix~\ref{app:monic-pivots}, together with that literal residual list,
specify every new polynomial by finite substitution.

\begin{theorem}[Acyclic monic elimination]
\label{thm:acyclic-monic}
Let $R_0,\ldots,R_{m-1}\in\Z[x_0,\ldots,x_{N-1}]$.
Select different rows and different coordinates so that each selected row
has the form
\[
 R_{r(v)}=\varepsilon_v x_v+f_v,
 \qquad \varepsilon_v\in\{-1,1\},
 \qquad x_v\text{ does not occur in }f_v.
\]
Suppose the directed graph on selected coordinates, with an edge $v\to w$
when $x_w$ occurs in $f_v$, is acyclic.  List the retained source indices
as $i_0<\cdots<i_{n-1}$, and let $y=(y_0,\ldots,y_{n-1})$.
Define integer polynomials $F_v(y)$ by $F_{i_j}=y_j$ on retained
coordinates, and recursively by
\[
 F_v=-\varepsilon_v f_v(F)
\]
on selected coordinates.  If $I$ is the set of unselected rows, put
\[
 Q(y)=\sum_{j\in I}R_j(F(y))^2.
\]
Projection to retained coordinates is a bijection between the integer
zero sets of $\sum_jR_j^2$ and $Q$.  Its inverse is the polynomial map
$y\mapsto(F_0(y),\ldots,F_{N-1}(y))$.
\end{theorem}

\begin{proof}
At an integer zero of the source sum of squares, every residual is zero.
Each selected row therefore forces exactly
$x_v=-\varepsilon_v f_v$.  Induction along a reverse topological order of the
acyclic graph gives $x_v=F_v(y)$ for every coordinate.  The unselected
residuals vanish, so $Q(y)=0$.

Conversely, if $Q(y)=0$, every remaining squared integer residual is zero.
Set $x=F(y)$.  The selected residuals vanish by their definitions, and
the unselected ones by $Q(y)=0$.  Thus the source polynomial is zero.
Every restored coordinate is an integer because all the substitutions
have integer coefficients.  The same induction proves uniqueness.
No division, nonzero denominator, congruence, or sign choice is needed.
\end{proof}

The unrestricted integer domain matters.  Applying this theorem to a
nonnegative-domain equation would require additional proofs that all
restored coordinates remain nonnegative at every zero.  We use the
integer source throughout.

\subsection{The resulting explicit family}

Every row in the following table defines one ordinary polynomial.  Its
complete arithmetic DAG uses only integer constants, its stated input
variables, and binary addition, subtraction and multiplication.
Intermediate registers abbreviate polynomials and are not additional
quantified unknowns.  Degree is exact.  The last column is a proved bound
on the height after collecting all coefficients; it is not the largest
constant appearing in the DAG.

\begin{center}
\begin{tabular}{rrrr}
\hline
Integer unknowns & Remaining residuals & Degree & Height bound \\
\hline
225 & 159 & 6 & $<2^{27}$ \\
153 & 87 & 8 & $<2^{47}$ \\
144 & 78 & 12 & $<2^{50}$ \\
125 & 59 & 16 & $<2^{80}$ \\
117 & 51 & 32 & $<2^{144}$ \\
111 & 45 & 64 & $<2^{173}$ \\
102 & 36 & 256 & $<2^{691}$ \\
94 & 28 & $65{,}540$ & $<2^{164870}$ \\
\hline
\end{tabular}
\end{center}

The degree-$16$ equation removes $123$ of the source's $248$ coordinates,
or $49.6\%$.  The degree-$256$ equation removes $146$, or $58.9\%$.
These are alternatives within the direct proof-calculus route, not
global dimension records among all coefficient encodings.  The
$94$-variable equation illustrates the cost of prioritizing dimension
alone: its degree and coefficient bound are much larger.

For the $153$-variable octic the package also supplies every collected
coefficient.  Its exact expansion has $188{,}579$ nonzero monomials and
height
\[
 750{,}273{,}658{,}864.
\]
Its $87$ fully substituted residuals contain $3{,}476$ monomials in total.
This exact height is stronger than the general $2^{47}$ bound.  The
expansion makes the support cost of eliminating coordinates explicit.

\subsection{Exact degree and coefficient bounds}

Give each retained coordinate degree one.  For each restored coordinate
propagate the maximum, over its monomials, of the sum of the previously
assigned coordinate degrees.  Apply the same rule to each remaining
residual.  If the largest bound is $d$, then $\deg Q\le2d$.

To prove equality, the supplied certificate gives a remaining residual
of degree bound $d$ whose degree-$d$ homogeneous part is nonzero when
the retained coordinate with source index $i$ is set to $i+3$, reduced
modulo $1{,}000{,}000{,}007$.  A nonzero modular value proves a nonzero
integer value; no primality assumption is needed.  At this integer point
the sum of the squares of all degree-$d$ homogeneous parts is positive.
It follows that $\deg Q=2d$.  This is an exact nonvanishing certificate,
not probabilistic polynomial identity testing.

\begin{lemma}[A propagated height bound]
\label{lem:monic-height}
Set $L_v=1$ for each retained coordinate.  For a restored polynomial
$F_v=\sum_M c_M\prod_{j\in M}F_j$, set
\[
 L_v=\sum_M |c_M|\prod_{j\in M}L_j.
\]
Use the same rule to obtain a bound $L_j'$ for each remaining residual.
Then every coefficient of $Q$ has absolute value at most
\[
 B_Q=\sum_{j\in I}(L_j')^2.
\]
If this positive integer has bit length $b$, every coefficient has
absolute value strictly below $2^b$.
\end{lemma}

\begin{proof}
For an integer polynomial $f$, let $\|f\|_1$ be the sum of the absolute
values of its collected coefficients.  The triangle and convolution
inequalities give
\[
 \|f+g\|_1\le\|f\|_1+\|g\|_1,
 \qquad \|fg\|_1\le\|f\|_1\|g\|_1.
\]
Induction through the substitutions proves $\|F_v\|_1\le L_v$.
Applying the same inequalities to the residuals and their squared sum
proves $\|Q\|_1\le B_Q$.  Each coefficient is bounded by this norm,
and $B_Q<2^b$ by the definition of bit length.
\end{proof}

All stated bit bounds are obtained by exact integer arithmetic.  The
largest literal constant in these DAGs is $256$; that different statistic
does not replace the expanded coefficient bounds in the table.

\subsection{Verification and the consistency theorem}

The independent checker expands each circuit block as a polynomial in
the coordinates already available.  These local expansions have
degree at most two, even when the fully substituted polynomial has degree
$65{,}540$.  It checks the blocks coefficient by coefficient against their
source rows, verifies the complete row partition and acyclic restoration,
and checks the final sum of squares.  A separate sparse engine expands
the octic completely.  Exact nonzero partial derivatives modulo
$1{,}000{,}000{,}007$ certify that every declared input occurs
in each collected polynomial.  These tests accompany the ordinary proof;
they do not replace its unbounded source calibration.

By Theorem~\ref{thm:acyclic-monic}, each displayed $Q$ has an integer zero
if and only if $P_{248}$ has one.  Compose this elementary, ZFC-formalizable
bijection with the already proved existence equivalence for the fused
front-sign source and its retained C222/ZFC calibration.  For every row
of the table this gives
\[
 \mathrm{ZFC}\ \vdash
 \bigl[\text{$Q$ has no zero in its stated integer domain}\bigr]
 \ \longleftrightarrow\ \operatorname{Con}(\mathrm{ZFC}).
\]
This establishes an equivalence, without assuming or proving either
side.  The new finite algebra is not claimed to be a fresh end-to-end
proof-assistant development.
% END elimination.tex


\section{Verification, reproduction, and use of the archive}
\label{sec:verification}

The archive is designed for local reproduction without a network connection.
Its mathematical data and checkers use Python's standard library.  The PDF
can be rebuilt with an ordinary pdfLaTeX installation.  From the package root:
\begin{verbatim}
python3 -B verify.py
python3 -B -O verify.py
python3 -B verify.py --rebuild
python3 build_report.py
\end{verbatim}
The first command checks the sealed file manifest and component/source
bindings, then runs the independent finite audits.  The second also disables
Python assertions throughout the component checks; no substantive check is
allowed to disappear with optimization.  The third additionally reconstructs
the delivered equation files and compares their bytes with the release.
The report builder updates the PDF and its single-file TeX edition.  Such a
rebuild intentionally changes a sealed document if the local TeX engine
differs, so a release's manifest should be checked before rebuilding it.

For the compact pair the separate checker binds the actual circuit gates
to $37$ exact polynomial cuts per equation, propagates the entire highest
homogeneous form, and verifies five additional algebraic identities.
It rejects $175$ actual gate mutations for the principal and $173$ for its
control.  The four fixed numerals are reconstructed from the complete seed
recipe, and every delivered coefficient payload is authenticated.  The
finite recurrence and index tests are diagnostic controls in addition to
the unbounded proof, not evidence that a consistency statement has been
decided by finite search.

For the direct quartics the checker independently reconstructs the changed
source, checks the $18$ Horn factoring identities, all monic projections,
all integral translations, the entire collected quartics, and the unchanged
parent file pins.  The principal has $203$ monic projection steps and $70$
translated coordinates.  A clean rebuild reproduces all $22$ generated
JSON artifacts for the three profiles.  Five deliberately corrupted new
artifacts are rejected; the parent's own audit is retained and run as well.

For the eight higher-degree forms the checker locally expands every circuit
block in previously defined source coordinates.  Every source row is
accounted for exactly once, either as a solved monic definition or as a
remaining residual.  This verifies the entire polynomial substitution
without expanding the highest-degree final forms.  Separate nonzero modular
derivative witnesses prove that every listed input actually occurs in the
resulting polynomial.  The $153$-variable octic is also expanded in full by a
separate sparse-polynomial engine.  Seven deliberate corruptions test the
algebraic certificate rejection path.  The complete coefficient norm bound
and a nonzero highest-degree witness establish the reported height bounds
and exact degrees; they are not floating-point estimates.

\subsection{Files that define the result}

The authoritative compact circuit is
\path{compact/equation_bV.json}.  The principal direct equation is
\path{direct_quartics/equations/integer248/equation.json}, together with its
fully collected \path{expanded_polynomial.json}.  The eight projected
circuits are in \path{higher_degree/equations/}.  Their source file is byte
identical to the principal $248$-variable equation; the top-level verifier
checks this binding explicitly.

All $182$ principal residuals are printed in Appendix~\ref{app:residuals}.
Together with the recursive elimination rule and the complete pivot tables
in Appendix~\ref{app:monic-pivots}, they give a finite paper specification of every
projected polynomial.  The JSON circuits are convenient executable versions
of those fixed definitions.  The complete fixed coefficient recipe for the
ten-positive form is printed in Appendix~\ref{app:constants}.

Each component contains its proof, producer, independent verifier, exact
evaluator or circuit format, source hashes and successful verification
receipts.  \path{review/} holds the additional scoped mathematical reviews
and root-level replays.  \path{SOURCE_LEDGER.json} records all inspected
incoming archive hashes.  \path{MANIFEST_SHA256.json} binds the final release
files.  Original licenses and notices are preserved.

\section{What the new ideas change, and what remains open}
\label{sec:outlook}

The compact result demonstrates a useful design principle: an index error
need not be encoded only to first order.  When a Pell/Lucas congruence falls
short of the modulus needed for a bounded rank, subtracting a short Taylor
term can increase usable precision without adding an existential variable.
Here the correction is cubic in the target rank, but its degree is far below
the dominant transport term.  The cost is a new signed-transport proof and
careful treatment of primes dividing its denominator.  The result succeeds
because those costs can be paid without destroying the final degree gain.

The direct result illustrates a different saving.  Two nonnegativity certificates
can sometimes be replaced by one product certificate if an existing exact
digit condition chooses the correct sign branch.  The difficult part is
the order of the proof: size bounds and power semantics must remain valid
in the unwanted branch, before factoriality or the desired signs are known.
The new signed-base argument supplies that order.

The higher-degree family uses standard monic substitution.  Its contribution
is the concrete, fully specified set of eliminations and certified tradeoff
points on the new source.  It should not be described as a new general
elimination theorem.  Compared with a quartic-only optimization, it exposes
how many of the remaining unknowns are merely arithmetic intermediate
values.  Its large-degree endpoint is included as a measured boundary point,
not selected as the most practical presentation.

Further reductions remain possible.  A still smaller substitution such as
$Y=q+s$ would change the crucial coprimality $H\equiv-1\pmod q$ and the
index-recovery divisibility; no proof for that substitution is supplied.
Higher-order corrections also need to pay for their signed tails and their
small-prime denominators.  For the direct family, the delivered projection
orders are achieved points, not a search-completeness certificate.  More
systematic degree/height-aware substitution or another semantic guard
recovery could improve them.  None of these prospects is counted as a
completed result in the tables.

\appendix
\clearpage
% BEGIN constants.tex
\section{Complete fixed coefficient recipe for the compact equation}
\label{app:constants}
This appendix reproduces the retained directly indexed BFS943 constant recipe.
Its local symbols are scoped to this appendix: the integer called $K$ here is
$K_0$ in Section~\ref{sec:lucas-q1}, and the radix $B$ here is unrelated to
the variable expression $B=(K_0+w)C-F$ there.  Every range is fixed.  Thus
this is a complete finite definition of four integer coefficients, not an
extra quantified condition.  The delivered binary coefficients are checked
against this recipe, including every digit of the seed.

The native tuple is inherited, with its state/color conjugacy and calibration;
it is not the new contribution of this paper.  See
\texttt{compact/parent/proof\_sources/indexed\_native.md} and the accompanying
source ledger for the precise boundary.

\textbf{The four fixed integers.}
Every index, set, sum, product and floor below defines a constant, not
another unknown. Put
\[
\rho=2^{67},\quad B=\rho^{8807},\quad\beta=B-1,\quad
G_m(x)=\sum_{j=0}^{m-1}x^j.
\]
Define
\[
\begin{split}
\mathcal E={}&\rg0{942}\cup\rg{983}{992}\cup\rg{1003}{1011}\\
&\cup\rg{1013}{1031}\cup\{1063,1064,1095,2195,3296\}.
\end{split}
\]
Then
\begin{align*}
M_C={}&B-1-\sum_{e\in\mathcal E\setminus\{1\}}\rho^e-2\rho^{1063},\\
M_F={}&B-1+4+(2^{51}-2)(\rho^{4393}+\rho^{5495})\\
&+\rho^{1013}G_{10}(\rho)+\rho^{1065}G_9(\rho)
  +\rho^{2195}+\rho^{3296}\\
&+(2^{65}-1)\rho^{8792}G_{13}(\rho)+(2^{19}-1)\rho^{8805},\\
K={}&\rho^{30}+\rho^{62}+\rho^{1100}+\rho^{2201}+\rho^{8849}\\
&+(\rho^{4393}+\rho^{5495})\Phi(1/\rho),
\end{align*}
where
\[
\begin{split}
\Phi(x)={}&(1+2^{51})G_{943}(x)
 -2^{10}\sum_{i=0}^{942}v(t_i)x^i\\
&+\sum_{(b,e,m)\in\mathcal C}2^b x^e G_m(2x),\\
\mathcal C={}&\{(10,1023,8),(18,1013,9),(27,1003,8),\\
&(35,983,9),(44,1031,1),(45,1022,1),(46,1011,1),\\
&(47,992,1),(48,1095,1),(49,2195,1),(50,3296,1)\}.
\end{split}
\]
The rational intermediate in K cancels exactly; all four results are positive integers.

\textbf{The 943 indexed fixed tiles.}
Use the actions $u_j,v_j,d_j$ decoded below. Let
\[
\begin{gathered}
h(q,b)=2(q+1)+b,\quad
\mathcal R=\{u_j:d_j=1,\ 0\le j<430\},\\
\mathcal L=\{u_j:d_j=0,\ 0\le j<430\},\\
r_q=6+|\mathcal R\cap\rg0{q-1}|,\quad
l_q=205+|\mathcal L\cap\rg0{q-1}|.
\end{gathered}
\]
An interval ending at $-1$ is empty. The table defines $t_i$ for
$0\le i\le942$; all index ranges are fixed:
\[
\setlength{\arraycolsep}{1.4pt}\begin{array}{c|l}
\text{index }i&\text{tile }t_i=(w,n,e,s)\\\hline
0,1&(2,435,3,1),\ (5,0,0,0)\\
j+2\ (j=2q+b<430)&((1-d_j)l_{u_j},h(q,b),d_jr_{u_j},v_j)\\
432+b\ (b=0,1)&(0,h(215,b),0,h(215,b))\\
434+b\ (b=0,1)&(0,b,0,b)\\
424+2r_q+b\ (q\in\mathcal R,b=0,1)&(r_q,b,0,h(q,b))\\
424+2l_q+b\ (q\in\mathcal L,b=0,1)&(0,b,l_q,h(q,b))\\
924+j\ (0\le j<4)&(0,j+430\fl{j/2},0,438)\\
928&(2,434,2,0)\\
929&(3,436,4,4)\\
930&(4,436,4,0)\\
497+t\ (t=434,435,436)&(1,438,1,t)\\
\substack{934+3a+b-6\fl{a/4}\\a\in\{0,1,4\},\ b=0,1,2}&(a,437,b,437)
\end{array}
\]
For $t=(w,n,e,s)$ put
\[
v(t)=w+2^8n+2^{17}e+2^{25}s+127\,2^{34}.
\]

\textbf{A closed decoder for the 430 actions.}
Put $D_b(A,k)=\fl{A/b^k}\bmod b$, where
$a\bmod b=a-b\fl{a/b}$. For $0\le q<215$, $b=0,1$, $j=2q+b$:
\[
\begin{gathered}
N=\mathcal A\bmod2^{636},\quad
T_0=\fl{\mathcal A/2^{636}}\bmod2^{1675},\\
E_0=\fl{\mathcal A/2^{2311}},\\
\nu_j=D_2(N,j),\quad \sigma_j=\sum_{k=0}^{j-1}\nu_k,\\
m_q=D_2(N,430+q),\quad \tau_q=\sum_{k=0}^{q-1}m_k,\\
F_q=14(1-m_q)+m_qD_{16}(E_0,\tau_q),\\
u_j=\nu_j(1+\sigma_j)+(1-\nu_j)D_{216}(T_0,j-\sigma_j),\\
v_j=D_2(F_q,2b),\quad d_j=D_2(F_q,2b+1).
\end{gathered}
\]
Empty sums are zero. Here $(v_0,d_0,u_0)=(1,1,1)$; the constructor tile
uses this actual first target, including its state renaming.

\textbf{The integer $\mathcal A$.}
The alphabet, in order of digit values $0,\ldots,61$, is\par
{\fontsize{10}{12}\selectfont\ttfamily 0123456789ABCDEFGHIJKLMNOPQRSTUVWXYZabcdefghijklmnopqrstuvwxyz}.\par
Concatenate the following 436 characters as digits $e_0,\ldots,e_{435}$;
put $\mathcal A=\sum_{j=0}^{435}e_j62^{435-j}$.\par
\begingroup\fontsize{10}{12}\selectfont\ttfamily
BxEnhvy6NipzkWmROo0GsA3uK3twWLxscLq20UTumy5Cdu461LwD6g7RDNzF8Y\par
iYDFTpB8ehAypd7d6OsOL4APamnJeNXQypUezCmW9n4oXP6XfgmiWuAWZ8ZJ7U\par
OP7gu4kVdlGZqZHULQMar1dXNeJW4i7oNVI8usp8BlovjjZbwPCovZI328fRUq\par
4q92R81s5GjWQaQfGhkq9UYMK3rRMaBmsOIBkJqc2cj2VRIZ6mjzkBIOkW9SjZ\par
FzeA6J5u4j3sBL3SdgwXzM9mHvQivu8EJ1TBTpRhHDqZIbmYgwQ8VqTAFotl6F\par
We1YC554GKrVHsA8oUGKSA6F6k5KREBSD0Tgd7bFALY84LLbpw574BQyvdoHrl\par
0Oe35uwDf15DpT6VJWZBjSbHPpPAINYTte0eUxnTFTKBlGj8Am6sb1gLldh8WZ\par
8P\par
\endgroup
% END constants.tex


\clearpage
\section{All residuals of the principal direct quartic}
\label{app:residuals}
The variables are exactly $x_0,\ldots,x_{247}\in\Z$.  Define
\[
 \boxed{\quad P_{248}(x)=\sum_{j=0}^{181}R_j(x)^2=0.\quad}
\]
The following list is complete and ordered.  There are no additional domain
restrictions or unlisted rows.  Its source file has SHA-256
\begin{center}\small
\path{c51ebc12534db042bd1c7bf7c59f84602f51625e2d952dffa82fa054b5d92192}.
\end{center}
The variable-to-source correspondence and every restoration are delivered
in the accompanying certificates.  For the higher-degree forms, the row
and coordinate numbers in this list are precisely those used by the pivot
tables; no renumbering is implicit.

\begingroup
\small
\setlength{\abovedisplayskip}{5pt}
\setlength{\belowdisplayskip}{5pt}
\setlength{\jot}{2pt}
% BEGIN residuals248.tex
% Residual 0: multiplication
\begin{align*}
R_{0} &= -x_{3}^{2} +x_{4}
\end{align*}

% Residual 1: multiplication
\begin{align*}
R_{1} &= -1 -x_{3} x_{5} +x_{5} +x_{6} -x_{7}
\end{align*}

% Residual 2: multiplication
\begin{align*}
R_{2} &= x_{0} -x_{0} x_{2} -x_{2} +x_{10}
\end{align*}

% Residual 3: multiplication
\begin{align*}
R_{3} &= -1 +x_{1} -x_{2} x_{8} -x_{3}
\end{align*}

% Residual 4: multiplication
\begin{align*}
R_{4} &= -1 +x_{1} -x_{9} x_{10}
\end{align*}

% Residual 5: multiplication
\begin{align*}
R_{5} &= 1 +x_{4} -x_{5} x_{21} +x_{19}
\end{align*}

% Residual 6: multiplication
\begin{align*}
R_{6} &= -1 +2 x_{11} -x_{11}^{2} -2 x_{11} x_{12} +2 x_{12} -x_{12}^{2} +x_{22}
\end{align*}

% Residual 7: multiplication
\begin{align*}
R_{7} &= -5 -30 x_{11} -36 x_{11}^{2} -96 x_{11} x_{12} -24 x_{11} x_{22} -40 x_{12} -64 x_{12}^{2} \\
&\quad -32 x_{12} x_{22} -10 x_{22} -4 x_{22}^{2} +x_{23}
\end{align*}

% Residual 8: multiplication
\begin{align*}
R_{8} &= -2 +2 x_{11} -x_{11}^{2} -4 x_{11} x_{23} +4 x_{23} -4 x_{23}^{2} +x_{24}
\end{align*}

% Residual 9: multiplication
\begin{align*}
R_{9} &= 1 +6 x_{11} -36 x_{11}^{2} -192 x_{11} x_{23} -24 x_{11} x_{24} +16 x_{23} -256 x_{23}^{2} \\
&\quad -64 x_{23} x_{24} +2 x_{24} -4 x_{24}^{2} +x_{25}
\end{align*}

% Residual 10: multiplication
\begin{align*}
R_{10} &= -2 -x_{12}^{2} -2 x_{12} x_{13} -x_{13}^{2} +x_{26}
\end{align*}

% Residual 11: multiplication
\begin{align*}
R_{11} &= -9 +7 x_{11} -x_{11}^{2} -4 x_{11} x_{12} -4 x_{11} x_{26} +12 x_{12} -4 x_{12}^{2} \\
&\quad -8 x_{12} x_{26} +12 x_{26} -4 x_{26}^{2} +x_{27}
\end{align*}

% Residual 12: multiplication
\begin{align*}
R_{12} &= -1 +2 x_{11} -x_{11}^{2} -2 x_{11} x_{13} +2 x_{13} -x_{13}^{2} +x_{28}
\end{align*}

% Residual 13: multiplication
\begin{align*}
R_{13} &= -6 x_{11} -144 x_{11}^{2} -192 x_{11} x_{12} -192 x_{11} x_{13} -48 x_{11} x_{22} \\
&\quad -48 x_{11} x_{28} -8 x_{12} -64 x_{12}^{2} -128 x_{12} x_{13} -32 x_{12} x_{22} -32 x_{12} x_{28} \\
&\quad -64 x_{13}^{2} -32 x_{13} x_{22} -32 x_{13} x_{28} -2 x_{22} -4 x_{22}^{2} -8 x_{22} x_{28} \\
&\quad -4 x_{28}^{2} +x_{29}
\end{align*}

% Residual 14: multiplication
\begin{align*}
R_{14} &= -4 x_{27}^{2} -8 x_{27} x_{29} -4 x_{29}^{2} +x_{30}
\end{align*}

% Residual 15: multiplication
\begin{align*}
R_{15} &= -6 +16 x_{13} -16 x_{13}^{2} -16 x_{13} x_{26} -16 x_{13} x_{28} +8 x_{26} -4 x_{26}^{2} \\
&\quad -8 x_{26} x_{28} +8 x_{28} -4 x_{28}^{2} +x_{31}
\end{align*}

% Residual 16: multiplication
\begin{align*}
R_{16} &= 1 +6 x_{11} -36 x_{11}^{2} -96 x_{11} x_{12} -48 x_{11} x_{13} -24 x_{11} x_{22} -48 x_{11} x_{28} \\
&\quad -24 x_{11} x_{31} +8 x_{12} -64 x_{12}^{2} -64 x_{12} x_{13} -32 x_{12} x_{22} -64 x_{12} x_{28} \\
&\quad -32 x_{12} x_{31} +8 x_{13} -16 x_{13}^{2} -16 x_{13} x_{22} -32 x_{13} x_{28} -16 x_{13} x_{31} \\
&\quad +2 x_{22} -4 x_{22}^{2} -16 x_{22} x_{28} -8 x_{22} x_{31} +8 x_{28} -16 x_{28}^{2} \\
&\quad -16 x_{28} x_{31} +4 x_{31} -4 x_{31}^{2} +x_{32}
\end{align*}

% Residual 17: multiplication
\begin{align*}
R_{17} &= -10 -5 x_{11} -x_{11}^{2} -12 x_{11} x_{12} -16 x_{11} x_{13} -4 x_{11} x_{26} -36 x_{12} \\
&\quad -36 x_{12}^{2} -96 x_{12} x_{13} -24 x_{12} x_{26} -48 x_{13} -64 x_{13}^{2} -32 x_{13} x_{26} \\
&\quad -12 x_{26} -4 x_{26}^{2} +x_{33}
\end{align*}

% Residual 18: multiplication
\begin{align*}
R_{18} &= -2 -24 x_{11} -36 x_{11}^{2} -12 x_{11} x_{12} -96 x_{11} x_{13} -24 x_{11} x_{28} -3 x_{12} \\
&\quad -x_{12}^{2} -16 x_{12} x_{13} -4 x_{12} x_{28} -32 x_{13} -64 x_{13}^{2} -32 x_{13} x_{28} \\
&\quad -8 x_{28} -4 x_{28}^{2} +x_{34}
\end{align*}

% Residual 19: multiplication
\begin{align*}
R_{19} &= 2 -4 x_{33}^{2} -8 x_{33} x_{34} -4 x_{34}^{2} +x_{35}
\end{align*}

% Residual 20: multiplication
\begin{align*}
R_{20} &= -12 +16 x_{11} -4 x_{11}^{2} -8 x_{11} x_{13} -8 x_{11} x_{22} -8 x_{11} x_{26} +16 x_{13} \\
&\quad -4 x_{13}^{2} -8 x_{13} x_{22} -8 x_{13} x_{26} +16 x_{22} -4 x_{22}^{2} -8 x_{22} x_{26} \\
&\quad +16 x_{26} -4 x_{26}^{2} +x_{36}
\end{align*}

% Residual 21: multiplication
\begin{align*}
R_{21} &= -3 -24 x_{11} -64 x_{11}^{2} -160 x_{11} x_{13} -64 x_{11} x_{26} -32 x_{11} x_{28} \\
&\quad -32 x_{11} x_{36} -26 x_{13} -100 x_{13}^{2} -80 x_{13} x_{26} -40 x_{13} x_{28} \\
&\quad -40 x_{13} x_{36} -8 x_{26} -16 x_{26}^{2} -16 x_{26} x_{28} -16 x_{26} x_{36} -6 x_{28} \\
&\quad -4 x_{28}^{2} -8 x_{28} x_{36} -4 x_{36} -4 x_{36}^{2} +x_{37}
\end{align*}

% Residual 22: multiplication
\begin{align*}
R_{22} &= -1 -18 x_{11} -36 x_{11}^{2} -96 x_{11} x_{12} -24 x_{11} x_{22} -24 x_{12} -64 x_{12}^{2} \\
&\quad -32 x_{12} x_{22} -6 x_{22} -4 x_{22}^{2} +x_{38}
\end{align*}

% Residual 23: multiplication
\begin{align*}
R_{23} &= -2 +2 x_{11} -x_{11}^{2} -4 x_{11} x_{38} +4 x_{38} -4 x_{38}^{2} +x_{39}
\end{align*}

% Residual 24: multiplication
\begin{align*}
R_{24} &= 1 +6 x_{11} -36 x_{11}^{2} -192 x_{11} x_{38} -24 x_{11} x_{39} +16 x_{38} -256 x_{38}^{2} \\
&\quad -64 x_{38} x_{39} +2 x_{39} -4 x_{39}^{2} +x_{40}
\end{align*}

% Residual 25: multiplication
\begin{align*}
R_{25} &= -x_{12}^{2} -4 x_{12} x_{40} -4 x_{40}^{2} +x_{41}
\end{align*}

% Residual 26: multiplication
\begin{align*}
R_{26} &= -6 -30 x_{11} -64 x_{11}^{2} -128 x_{11} x_{12} -64 x_{11} x_{22} -32 x_{12} -64 x_{12}^{2} \\
&\quad -64 x_{12} x_{22} -14 x_{22} -16 x_{22}^{2} +x_{42}
\end{align*}

% Residual 27: multiplication
\begin{align*}
R_{27} &= -4 +8 x_{12} -64 x_{12}^{2} -96 x_{12} x_{13} -32 x_{12} x_{26} -32 x_{12} x_{42} +6 x_{13} \\
&\quad -36 x_{13}^{2} -24 x_{13} x_{26} -24 x_{13} x_{42} +2 x_{26} -4 x_{26}^{2} -8 x_{26} x_{42} \\
&\quad +4 x_{42} -4 x_{42}^{2} +x_{43}
\end{align*}

% Residual 28: multiplication
\begin{align*}
R_{28} &= 8 x_{11} -64 x_{11}^{2} -96 x_{11} x_{13} -32 x_{11} x_{28} -32 x_{11} x_{43} +6 x_{13} \\
&\quad -36 x_{13}^{2} -24 x_{13} x_{28} -24 x_{13} x_{43} +2 x_{28} -4 x_{28}^{2} -8 x_{28} x_{43} \\
&\quad +4 x_{43} -4 x_{43}^{2} +x_{44}
\end{align*}

% Residual 29: multiplication
\begin{align*}
R_{29} &= -x_{13}^{2} -2 x_{13} x_{14} -x_{14}^{2} +x_{45}
\end{align*}

% Residual 30: multiplication
\begin{align*}
R_{30} &= 2 x_{11} -x_{11}^{2} -4 x_{11} x_{13} -4 x_{11} x_{45} +4 x_{13} -4 x_{13}^{2} -8 x_{13} x_{45} \\
&\quad +4 x_{45} -4 x_{45}^{2} +x_{46}
\end{align*}

% Residual 31: multiplication
\begin{align*}
R_{31} &= -35 +24 x_{11} -4 x_{11}^{2} -8 x_{11} x_{12} -8 x_{11} x_{26} -8 x_{11} x_{28} +24 x_{12} \\
&\quad -4 x_{12}^{2} -8 x_{12} x_{26} -8 x_{12} x_{28} +24 x_{26} -4 x_{26}^{2} -8 x_{26} x_{28} \\
&\quad +24 x_{28} -4 x_{28}^{2} +x_{47}
\end{align*}

% Residual 32: multiplication
\begin{align*}
R_{32} &= -16 x_{11}^{2} -32 x_{11} x_{13} -32 x_{11} x_{28} -32 x_{11} x_{45} -16 x_{11} x_{46} \\
&\quad -16 x_{11} x_{47} -4 x_{13} -16 x_{13}^{2} -32 x_{13} x_{28} -32 x_{13} x_{45} -16 x_{13} x_{46} \\
&\quad -16 x_{13} x_{47} -16 x_{28}^{2} -32 x_{28} x_{45} -16 x_{28} x_{46} -16 x_{28} x_{47} -4 x_{45} \\
&\quad -16 x_{45}^{2} -16 x_{45} x_{46} -16 x_{45} x_{47} -2 x_{46} -4 x_{46}^{2} -8 x_{46} x_{47} \\
&\quad -4 x_{47}^{2} +x_{48}
\end{align*}

% Residual 33: multiplication
\begin{align*}
R_{33} &= x_{13} -25 x_{13}^{2} -40 x_{13} x_{28} -20 x_{13} x_{31} -16 x_{28}^{2} -16 x_{28} x_{31} \\
&\quad -4 x_{31}^{2} +x_{49}
\end{align*}

% Residual 34: multiplication
\begin{align*}
R_{34} &= -1 +8 x_{12} -4 x_{12}^{2} -8 x_{12} x_{26} +8 x_{26} -4 x_{26}^{2} +x_{50}
\end{align*}

% Residual 35: multiplication
\begin{align*}
R_{35} &= -1 -16 x_{12}^{2} -8 x_{12} x_{13} -32 x_{12} x_{26} -16 x_{12} x_{50} -x_{13}^{2} \\
&\quad -8 x_{13} x_{26} -4 x_{13} x_{50} -16 x_{26}^{2} -16 x_{26} x_{50} -4 x_{50}^{2} +x_{51}
\end{align*}

% Residual 36: multiplication
\begin{align*}
R_{36} &= 1 +x_{12} -x_{12}^{2} -4 x_{12} x_{13} -4 x_{12} x_{51} -4 x_{13}^{2} -8 x_{13} x_{51} \\
&\quad -4 x_{51}^{2} +x_{52}
\end{align*}

% Residual 37: multiplication
\begin{align*}
R_{37} &= -1 -4 x_{52} -4 x_{52}^{2} +x_{53}
\end{align*}

% Residual 38: multiplication
\begin{align*}
R_{38} &= -1 -x_{13}^{2} -8 x_{13} x_{52} -4 x_{13} x_{53} -16 x_{52}^{2} -16 x_{52} x_{53} -4 x_{53}^{2} \\
&\quad +x_{54}
\end{align*}

% Residual 39: multiplication
\begin{align*}
R_{39} &= 2 +x_{12} -x_{12}^{2} -4 x_{12} x_{13} -4 x_{12} x_{54} -4 x_{13}^{2} -8 x_{13} x_{54} \\
&\quad -4 x_{54}^{2} +x_{55}
\end{align*}

% Residual 40: multiplication
\begin{align*}
R_{40} &= -4 x_{49}^{2} -8 x_{49} x_{55} -4 x_{55}^{2} +x_{56}
\end{align*}

% Residual 41: multiplication
\begin{align*}
R_{41} &= -9 +7 x_{11} -x_{11}^{2} -8 x_{11} x_{49} -4 x_{11} x_{56} +24 x_{49} -16 x_{49}^{2} \\
&\quad -16 x_{49} x_{56} +12 x_{56} -4 x_{56}^{2} +x_{57}
\end{align*}

% Residual 42: multiplication
\begin{align*}
R_{42} &= -1 -4 x_{57} -4 x_{57}^{2} +x_{58}
\end{align*}

% Residual 43: multiplication
\begin{align*}
R_{43} &= -x_{13}^{2} -8 x_{13} x_{57} -4 x_{13} x_{58} -16 x_{57}^{2} -16 x_{57} x_{58} -4 x_{58}^{2} \\
&\quad +x_{59}
\end{align*}

% Residual 44: multiplication
\begin{align*}
R_{44} &= -x_{11} -x_{11}^{2} -4 x_{11} x_{13} -4 x_{11} x_{59} -4 x_{13} -4 x_{13}^{2} -8 x_{13} x_{59} \\
&\quad -4 x_{59} -4 x_{59}^{2} +x_{60}
\end{align*}

% Residual 45: multiplication
\begin{align*}
R_{45} &= -10 -12 x_{60} -4 x_{60}^{2} +x_{61}
\end{align*}

% Residual 46: multiplication
\begin{align*}
R_{46} &= 1 +4 x_{12} -4 x_{12}^{2} -8 x_{12} x_{26} -16 x_{12} x_{60} -8 x_{12} x_{61} +4 x_{26} \\
&\quad -4 x_{26}^{2} -16 x_{26} x_{60} -8 x_{26} x_{61} +8 x_{60} -16 x_{60}^{2} -16 x_{60} x_{61} \\
&\quad +4 x_{61} -4 x_{61}^{2} +x_{62}
\end{align*}

% Residual 47: multiplication
\begin{align*}
R_{47} &= x_{12} -25 x_{12}^{2} -40 x_{12} x_{26} -20 x_{12} x_{62} -16 x_{26}^{2} -16 x_{26} x_{62} \\
&\quad -4 x_{62}^{2} +x_{63}
\end{align*}

% Residual 48: multiplication
\begin{align*}
R_{48} &= -10 -12 x_{55} -4 x_{55}^{2} +x_{64}
\end{align*}

% Residual 49: multiplication
\begin{align*}
R_{49} &= 1 -16 x_{55}^{2} -16 x_{55} x_{63} -16 x_{55} x_{64} -4 x_{63}^{2} -8 x_{63} x_{64} -4 x_{64}^{2} \\
&\quad +x_{65}
\end{align*}

% Residual 50: multiplication
\begin{align*}
R_{50} &= x_{13} -x_{13}^{2} -8 x_{13} x_{63} -4 x_{13} x_{65} -16 x_{63}^{2} -16 x_{63} x_{65} \\
&\quad -4 x_{65}^{2} +x_{66}
\end{align*}

% Residual 51: multiplication
\begin{align*}
R_{51} &= -4 x_{11} -4 x_{11}^{2} -4 x_{11} x_{13} -8 x_{11} x_{22} -x_{13} -x_{13}^{2} -4 x_{13} x_{22} \\
&\quad -4 x_{22} -4 x_{22}^{2} +x_{67}
\end{align*}

% Residual 52: multiplication
\begin{align*}
R_{52} &= 1 +4 x_{11} -4 x_{11}^{2} -8 x_{11} x_{22} -8 x_{11} x_{67} +4 x_{22} -4 x_{22}^{2} \\
&\quad -8 x_{22} x_{67} +4 x_{67} -4 x_{67}^{2} +x_{68}
\end{align*}

% Residual 53: multiplication
\begin{align*}
R_{53} &= -2 -4 x_{11} -16 x_{11}^{2} -32 x_{11} x_{22} -16 x_{11} x_{66} -16 x_{11} x_{68} -4 x_{22} \\
&\quad -16 x_{22}^{2} -16 x_{22} x_{66} -16 x_{22} x_{68} -4 x_{66}^{2} -8 x_{66} x_{68} -2 x_{68} \\
&\quad -4 x_{68}^{2} +x_{69}
\end{align*}

% Residual 54: multiplication
\begin{align*}
R_{54} &= -2 -4 x_{69} -4 x_{69}^{2} +x_{70}
\end{align*}

% Residual 55: multiplication
\begin{align*}
R_{55} &= -x_{3} x_{16} +x_{71}
\end{align*}

% Residual 56: multiplication
\begin{align*}
R_{56} &= x_{15} -x_{16} -2 x_{25} x_{71} -x_{72}
\end{align*}

% Residual 57: multiplication
\begin{align*}
R_{57} &= x_{15} -x_{16} -4 x_{27} x_{71} -2 x_{30} x_{71} -x_{73}
\end{align*}

% Residual 58: multiplication
\begin{align*}
R_{58} &= x_{15} -x_{16} -2 x_{32} x_{71} -6 x_{71} -x_{74}
\end{align*}

% Residual 59: multiplication
\begin{align*}
R_{59} &= x_{15} -x_{16} -4 x_{33} x_{71} -2 x_{35} x_{71} -x_{75}
\end{align*}

% Residual 60: multiplication
\begin{align*}
R_{60} &= x_{15} -x_{16} -2 x_{37} x_{71} +6 x_{71} -x_{76}
\end{align*}

% Residual 61: multiplication
\begin{align*}
R_{61} &= x_{15} -x_{16} -4 x_{40} x_{71} -2 x_{41} x_{71} -x_{77}
\end{align*}

% Residual 62: multiplication
\begin{align*}
R_{62} &= x_{15} -x_{16} -2 x_{44} x_{71} -x_{78}
\end{align*}

% Residual 63: multiplication
\begin{align*}
R_{63} &= x_{15} -x_{16} -2 x_{48} x_{71} -4 x_{71} -x_{79}
\end{align*}

% Residual 64: nonnegative_square_difference
\begin{align*}
R_{64} &= -x_{11} x_{12} +x_{11} x_{13} +x_{12} +x_{12} x_{13} -x_{13} -x_{13}^{2} +x_{80}
\end{align*}

% Residual 65: multiplication
\begin{align*}
R_{65} &= 1 +2 x_{14} +x_{14}^{2} -x_{80}^{2} +x_{81}
\end{align*}

% Residual 66: nonnegative_square_difference
\begin{align*}
R_{66} &= -x_{81}^{2} +x_{82}
\end{align*}

% Residual 67: multiplication
\begin{align*}
R_{67} &= -4 x_{69} x_{71} -2 x_{70} x_{71} +2 x_{71} +x_{83}
\end{align*}

% Residual 68: nonnegative_square_difference
\begin{align*}
R_{68} &= -x_{15}^{2} +2 x_{15} x_{16} +2 x_{15} x_{83} -x_{16}^{2} -2 x_{16} x_{83} -x_{83}^{2} +x_{84}
\end{align*}

% Residual 69: multiplication
\begin{align*}
R_{69} &= -1 -x_{3} x_{12} +x_{85}
\end{align*}

% Residual 70: multiplication
\begin{align*}
R_{70} &= -x_{16} x_{85} +x_{87} x_{88} +x_{89}
\end{align*}

% Residual 71: multiplication
\begin{align*}
R_{71} &= -1 +x_{3} -x_{3} x_{11} +x_{86}
\end{align*}

% Residual 72: multiplication
\begin{align*}
R_{72} &= -1 +2 x_{3} -2 x_{3} x_{11} -2 x_{3} x_{22} +x_{87}
\end{align*}

% Residual 73: multiplication
\begin{align*}
R_{73} &= -x_{15} x_{86} +x_{88}
\end{align*}

% Residual 74: multiplication
\begin{align*}
R_{74} &= -1 -6 x_{3} x_{11} -8 x_{3} x_{12} -2 x_{3} x_{22} +x_{90}
\end{align*}

% Residual 75: multiplication
\begin{align*}
R_{75} &= x_{15} x_{85} -x_{16} x_{90} +x_{91}
\end{align*}

% Residual 76: multiplication
\begin{align*}
R_{76} &= -x_{15} x_{72} +x_{16} x_{72} +x_{92}
\end{align*}

% Residual 77: multiplication
\begin{align*}
R_{77} &= -x_{73} x_{92} +x_{93}
\end{align*}

% Residual 78: multiplication
\begin{align*}
R_{78} &= -x_{74} x_{93} +x_{94}
\end{align*}

% Residual 79: multiplication
\begin{align*}
R_{79} &= -x_{75} x_{94} +x_{95}
\end{align*}

% Residual 80: multiplication
\begin{align*}
R_{80} &= -x_{76} x_{95} +x_{96}
\end{align*}

% Residual 81: multiplication
\begin{align*}
R_{81} &= -x_{77} x_{96} +x_{97}
\end{align*}

% Residual 82: multiplication
\begin{align*}
R_{82} &= -x_{78} x_{97} +x_{98}
\end{align*}

% Residual 83: multiplication
\begin{align*}
R_{83} &= -x_{79} x_{98} +x_{99}
\end{align*}

% Residual 84: multiplication
\begin{align*}
R_{84} &= -x_{89} x_{99} +x_{100}
\end{align*}

% Residual 85: multiplication
\begin{align*}
R_{85} &= -x_{91} x_{100} +x_{101}
\end{align*}

% Residual 86: multiplication
\begin{align*}
R_{86} &= -x_{101}^{2} +x_{102}
\end{align*}

% Residual 87: multiplication
\begin{align*}
R_{87} &= -x_{2} -x_{2} x_{19} +x_{19} +x_{103}
\end{align*}

% Residual 88: multiplication
\begin{align*}
R_{88} &= -x_{17} x_{103} +x_{104}
\end{align*}

% Residual 89: nonnegative_square_difference
\begin{align*}
R_{89} &= -2 x_{1}^{2} +2 x_{1} x_{15} +2 x_{1} x_{16} +2 x_{1} x_{104} +2 x_{1} x_{105} -x_{15}^{2} \\
&\quad -2 x_{15} x_{104} -x_{16}^{2} -2 x_{16} x_{105} +x_{20} x_{107} -x_{82} x_{102} -x_{84} x_{102} \\
&\quad -x_{104}^{2} -x_{105}^{2}
\end{align*}

% Residual 90: multiplication
\begin{align*}
R_{90} &= -x_{2} x_{18} +x_{18} -x_{18} x_{103} +x_{105}
\end{align*}

% Residual 91: addition
\begin{align*}
R_{91} &= -9 x_{11} -9 x_{12} -4 x_{22} -2 x_{25} -4 x_{27} -2 x_{30} -2 x_{32} -4 x_{33} -2 x_{35} -2 x_{37} \\
&\quad -4 x_{40} -2 x_{41} -2 x_{44} -2 x_{48} -4 x_{69} -2 x_{70} +x_{106}
\end{align*}

% Residual 92: multiplication
\begin{align*}
R_{92} &= -x_{108} x_{111} +x_{111} +x_{113}
\end{align*}

% Residual 93: choose0.digit
\begin{align*}
R_{93} &= -x_{5} x_{135} +x_{109} x_{235} -x_{112} -x_{113} x_{135}
\end{align*}

% Residual 94: crt.range.a.quotient_multiple
\begin{align*}
R_{94} &= x_{114} x_{236} -x_{116} x_{138} -x_{117}
\end{align*}

% Residual 95: crt.range.b.quotient_multiple
\begin{align*}
R_{95} &= x_{118} x_{236} -x_{119} x_{138} -x_{120}
\end{align*}

% Residual 96: crt.range.c.quotient_multiple
\begin{align*}
R_{96} &= x_{121} x_{236} -x_{122} x_{138} -x_{123}
\end{align*}

% Residual 97: crt.range.d.quotient_multiple
\begin{align*}
R_{97} &= x_{124} x_{236} -x_{125} x_{138} -x_{126}
\end{align*}

% Residual 98: crt.range.S.quotient_multiple
\begin{align*}
R_{98} &= x_{127} x_{237} -x_{129} x_{138} -x_{130}
\end{align*}

% Residual 99: crt.range.Sp.quotient_multiple
\begin{align*}
R_{99} &= x_{131} x_{237} -x_{132} x_{138} -x_{133}
\end{align*}

% Residual 100: multiplication
\begin{align*}
R_{100} &= -x_{20} x_{135} +x_{138}
\end{align*}

% Residual 101: multiplication
\begin{align*}
R_{101} &= x_{134} -x_{136} x_{138} -x_{137}
\end{align*}

% Residual 102: multiplication
\begin{align*}
R_{102} &= -x_{138} x_{140} +x_{139} -x_{141}
\end{align*}

% Residual 103: crt.range.index.quotient_multiple
\begin{align*}
R_{103} &= -x_{138} x_{143} +x_{142} x_{235} -x_{144}
\end{align*}

% Residual 104: crt.range.derived_label_sum.quotient_multiple
\begin{align*}
R_{104} &= -x_{138} x_{146} +x_{145} x_{236} -x_{147}
\end{align*}

% Residual 105: multiplication
\begin{align*}
R_{105} &= x_{148} -2 x_{148} x_{150} +x_{151}
\end{align*}

% Residual 106: multiplication
\begin{align*}
R_{106} &= -x_{148} x_{151} -x_{151} +x_{152}
\end{align*}

% Residual 107: shared.last_norm
\begin{align*}
R_{107} &= 2 x_{149} +x_{149}^{2} +x_{151}^{2} -x_{152}^{2}
\end{align*}

% Residual 108: multiplication
\begin{align*}
R_{108} &= -x_{149} x_{154} -x_{154} +x_{155}
\end{align*}

% Residual 109: index0.norm1
\begin{align*}
R_{109} &= 2 x_{153} +x_{153}^{2} +x_{154}^{2} -x_{155}^{2}
\end{align*}

% Residual 110: multiplication
\begin{align*}
R_{110} &= x_{4}^{2} -2 x_{4} x_{149} -2 x_{149} +x_{156}
\end{align*}

% Residual 111: multiplication
\begin{align*}
R_{111} &= -1 -x_{4} x_{154} +x_{6} -x_{153} -x_{154} +x_{155} +x_{156} x_{157}
\end{align*}

% Residual 112: multiplication
\begin{align*}
R_{112} &= 1 -2 x_{5} +x_{5}^{2} -2 x_{5} x_{149} +x_{158}
\end{align*}

% Residual 113: multiplication
\begin{align*}
R_{113} &= -x_{5} x_{204} +x_{20} +x_{158} x_{159} -x_{203} +x_{205}
\end{align*}

% Residual 114: base2.T_definition
\begin{align*}
R_{114} &= -1 +2 x_{149} x_{219} -x_{160} +2 x_{219} -x_{219}^{2}
\end{align*}

% Residual 115: multiplication
\begin{align*}
R_{115} &= -1 +x_{108} +x_{160} x_{161} -x_{203} -x_{204} x_{219} +x_{205}
\end{align*}

% Residual 116: multiplication
\begin{align*}
R_{116} &= -x_{149} x_{163} -x_{163} +x_{164}
\end{align*}

% Residual 117: index2.norm1
\begin{align*}
R_{117} &= 2 x_{162} +x_{162}^{2} +x_{163}^{2} -x_{164}^{2}
\end{align*}

% Residual 118: multiplication
\begin{align*}
R_{118} &= 1 -2 x_{108} +x_{108}^{2} -2 x_{108} x_{149} +x_{165}
\end{align*}

% Residual 119: multiplication
\begin{align*}
R_{119} &= -1 -x_{108} x_{163} +x_{109} -x_{162} +x_{164} +x_{165} x_{166}
\end{align*}

% Residual 120: multiplication
\begin{align*}
R_{120} &= 4 -4 x_{108} +x_{108}^{2} -2 x_{108} x_{149} +2 x_{149} +x_{167}
\end{align*}

% Residual 121: multiplication
\begin{align*}
R_{121} &= -1 -x_{108} x_{154} +x_{110} -x_{153} +x_{154} +x_{155} +x_{167} x_{168}
\end{align*}

% Residual 122: multiplication
\begin{align*}
R_{122} &= 1 +x_{149} -x_{149} x_{170} -x_{170} +x_{171}
\end{align*}

% Residual 123: index3.norm1
\begin{align*}
R_{123} &= 1 +2 x_{169} +x_{169}^{2} -2 x_{170} +x_{170}^{2} -x_{171}^{2}
\end{align*}

% Residual 124: multiplication
\begin{align*}
R_{124} &= -1 +x_{108} -x_{108} x_{170} +x_{114} +x_{165} x_{172} -x_{169} +x_{171}
\end{align*}

% Residual 125: multiplication
\begin{align*}
R_{125} &= -2 x_{149} -x_{149}^{2} +x_{173}
\end{align*}

% Residual 126: multiplication
\begin{align*}
R_{126} &= 1 -x_{108} -x_{108} x_{153} -x_{108} x_{155} +x_{115} +x_{149} x_{155} +x_{153} -x_{154} x_{173} \\
&\quad +2 x_{155} +x_{167} x_{174}
\end{align*}

% Residual 127: multiplication
\begin{align*}
R_{127} &= -x_{149} x_{176} -x_{176} +x_{177}
\end{align*}

% Residual 128: index5.norm1
\begin{align*}
R_{128} &= 2 x_{175} +x_{175}^{2} +x_{176}^{2} -x_{177}^{2}
\end{align*}

% Residual 129: multiplication
\begin{align*}
R_{129} &= -1 -x_{108} x_{176} +x_{118} +x_{165} x_{178} -x_{175} +x_{177}
\end{align*}

% Residual 130: multiplication
\begin{align*}
R_{130} &= -x_{149} x_{180} -x_{180} +x_{181}
\end{align*}

% Residual 131: index6.norm1
\begin{align*}
R_{131} &= 2 x_{179} +x_{179}^{2} +x_{180}^{2} -x_{181}^{2}
\end{align*}

% Residual 132: multiplication
\begin{align*}
R_{132} &= -1 -x_{108} x_{180} +x_{121} +x_{165} x_{182} -x_{179} +x_{181}
\end{align*}

% Residual 133: multiplication
\begin{align*}
R_{133} &= -x_{149} x_{184} -x_{184} +x_{185}
\end{align*}

% Residual 134: index7.norm1
\begin{align*}
R_{134} &= 2 x_{183} +x_{183}^{2} +x_{184}^{2} -x_{185}^{2}
\end{align*}

% Residual 135: multiplication
\begin{align*}
R_{135} &= -1 -x_{108} x_{184} +x_{124} +x_{165} x_{186} -x_{183} +x_{185}
\end{align*}

% Residual 136: multiplication
\begin{align*}
R_{136} &= -x_{149} x_{188} -x_{188} +x_{189}
\end{align*}

% Residual 137: index8.norm1
\begin{align*}
R_{137} &= 2 x_{187} +x_{187}^{2} +x_{188}^{2} -x_{189}^{2}
\end{align*}

% Residual 138: multiplication
\begin{align*}
R_{138} &= -1 -x_{108} x_{188} +x_{127} +x_{165} x_{190} -x_{187} +x_{189}
\end{align*}

% Residual 139: multiplication
\begin{align*}
R_{139} &= -x_{149} x_{192} -x_{192} +x_{193}
\end{align*}

% Residual 140: index9.norm1
\begin{align*}
R_{140} &= 2 x_{191} +x_{191}^{2} +x_{192}^{2} -x_{193}^{2}
\end{align*}

% Residual 141: multiplication
\begin{align*}
R_{141} &= -1 -x_{108} x_{192} +x_{128} +x_{167} x_{194} -x_{191} +x_{192} +x_{193}
\end{align*}

% Residual 142: multiplication
\begin{align*}
R_{142} &= -x_{149} x_{196} -x_{196} +x_{197}
\end{align*}

% Residual 143: index10.norm1
\begin{align*}
R_{143} &= 2 x_{195} +x_{195}^{2} +x_{196}^{2} -x_{197}^{2}
\end{align*}

% Residual 144: multiplication
\begin{align*}
R_{144} &= -1 -x_{108} x_{196} +x_{131} +x_{165} x_{198} -x_{195} +x_{197}
\end{align*}

% Residual 145: multiplication
\begin{align*}
R_{145} &= -x_{149} x_{200} -x_{200} +x_{201}
\end{align*}

% Residual 146: index11.norm1
\begin{align*}
R_{146} &= 2 x_{199} +x_{199}^{2} +x_{200}^{2} -x_{201}^{2}
\end{align*}

% Residual 147: multiplication
\begin{align*}
R_{147} &= -1 -x_{108} x_{200} +x_{134} +x_{165} x_{202} -x_{199} +x_{201}
\end{align*}

% Residual 148: multiplication
\begin{align*}
R_{148} &= 1 -x_{149} x_{204} -x_{204} +x_{205}
\end{align*}

% Residual 149: index12.norm1
\begin{align*}
R_{149} &= -1 +2 x_{203} +x_{203}^{2} +x_{204}^{2} -2 x_{205} -x_{205}^{2}
\end{align*}

% Residual 150: multiplication
\begin{align*}
R_{150} &= -x_{108} x_{204} +x_{135} +x_{167} x_{206} -x_{203} +x_{204} +x_{205}
\end{align*}

% Residual 151: multiplication
\begin{align*}
R_{151} &= -x_{149} x_{208} -x_{208} +x_{209}
\end{align*}

% Residual 152: index13.norm1
\begin{align*}
R_{152} &= 2 x_{207} +x_{207}^{2} +x_{208}^{2} -x_{209}^{2}
\end{align*}

% Residual 153: multiplication
\begin{align*}
R_{153} &= -1 -x_{108} x_{208} +x_{139} +x_{165} x_{210} -x_{207} +x_{209}
\end{align*}

% Residual 154: multiplication
\begin{align*}
R_{154} &= -x_{149} x_{212} -x_{212} +x_{213}
\end{align*}

% Residual 155: index14.norm1
\begin{align*}
R_{155} &= 2 x_{211} +x_{211}^{2} +x_{212}^{2} -x_{213}^{2}
\end{align*}

% Residual 156: multiplication
\begin{align*}
R_{156} &= -1 -x_{108} x_{212} +x_{142} +x_{165} x_{214} -x_{211} +x_{213}
\end{align*}

% Residual 157: multiplication
\begin{align*}
R_{157} &= -x_{149} x_{216} -x_{216} +x_{217}
\end{align*}

% Residual 158: index15.norm1
\begin{align*}
R_{158} &= 2 x_{215} +x_{215}^{2} +x_{216}^{2} -x_{217}^{2}
\end{align*}

% Residual 159: multiplication
\begin{align*}
R_{159} &= -1 -x_{108} x_{216} +x_{145} +x_{165} x_{218} -x_{215} +x_{217}
\end{align*}

% Residual 160: shared.frame_definition
\begin{align*}
R_{160} &= -6 -2 x_{0} -3 x_{0}^{2} -2 x_{0} x_{4} -x_{1}^{2} -x_{2}^{2} -2 x_{4} -2 x_{4}^{2} -x_{5}^{2} \\
&\quad -x_{6}^{2} +2 x_{11} -x_{11}^{2} -x_{12}^{2} -x_{13}^{2} -x_{14}^{2} -x_{15}^{2} -x_{16}^{2} \\
&\quad -x_{17}^{2} -x_{18}^{2} -x_{19}^{2} -x_{20}^{2} -x_{106}^{2} +2 x_{108} -2 x_{108}^{2} -x_{109}^{2} \\
&\quad -x_{110}^{2} -x_{114}^{2} -x_{115}^{2} -x_{118}^{2} -x_{121}^{2} -x_{124}^{2} -x_{127}^{2} \\
&\quad -x_{128}^{2} -x_{131}^{2} -x_{134}^{2} -x_{135}^{2} -x_{139}^{2} -x_{142}^{2} -x_{145}^{2} +x_{148} \\
&\quad -x_{219}^{2}
\end{align*}

% Residual 161: self_radix.base_definition
\begin{align*}
R_{161} &= -x_{5} x_{11} -x_{5} x_{12} -x_{5} x_{13} -x_{5} x_{14} -x_{5} x_{15} -x_{5} x_{16} -x_{5} x_{17} \\
&\quad -x_{5} x_{18} -x_{5} x_{19} +x_{219}
\end{align*}

% Residual 162: cofinal.horizon
\begin{align*}
R_{162} &= -16 +x_{0} -x_{220}^{2}
\end{align*}

% Residual 163: cofinal.beta_base
\begin{align*}
R_{163} &= -1 +x_{2} +x_{6} -x_{6}^{2}
\end{align*}

% Residual 164: cap.common_Ay
\begin{align*}
R_{164} &= -x_{149} x_{222} -x_{222} +x_{231}
\end{align*}

% Residual 165: cap.common_norm
\begin{align*}
R_{165} &= -1 +x_{221}^{2} +x_{222}^{2} -x_{231}^{2}
\end{align*}

% Residual 166: cap.BR
\begin{align*}
R_{166} &= -x_{224} -x_{224} x_{225} +x_{232}
\end{align*}

% Residual 167: cap.comparison_norm
\begin{align*}
R_{167} &= -1 +x_{223}^{2} +x_{224}^{2} -x_{232}^{2}
\end{align*}

% Residual 168: cap.Y2
\begin{align*}
R_{168} &= -x_{216}^{2} +x_{233}
\end{align*}

% Residual 169: cap.common_multiple
\begin{align*}
R_{169} &= x_{222} -2 x_{227} x_{233} -x_{233}
\end{align*}

% Residual 170: cap.comparison_constructor
\begin{align*}
R_{170} &= -4 x_{222} -8 x_{222} x_{226} +x_{225}
\end{align*}

% Residual 171: cap.comparison_modulus
\begin{align*}
R_{171} &= -x_{149} -x_{221} x_{228} +x_{225}
\end{align*}

% Residual 172: cap.abscissa_residue
\begin{align*}
R_{172} &= -1 -x_{215} -x_{221} x_{229} +x_{223}
\end{align*}

% Residual 173: cap.index_residue
\begin{align*}
R_{173} &= -x_{106} -4 x_{216} x_{230} +x_{224}
\end{align*}

% Residual 174: cap.pooled_index_residue
\begin{align*}
R_{174} &= 2 x_{0}^{2} +2 x_{0} x_{4} -2 x_{0} x_{154} -2 x_{0} x_{163} +x_{1}^{2} -2 x_{1} x_{204} \\
&\quad +x_{2}^{2} -2 x_{2} x_{192} +x_{4}^{2} -2 x_{4} x_{163} +x_{11}^{2} -2 x_{11} x_{170} +x_{12}^{2} \\
&\quad -2 x_{12} x_{176} +x_{13}^{2} -2 x_{13} x_{180} +x_{14}^{2} -2 x_{14} x_{184} +x_{15}^{2} \\
&\quad -2 x_{15} x_{188} +x_{16}^{2} -2 x_{16} x_{196} +x_{17}^{2} -2 x_{17} x_{200} +x_{18}^{2} \\
&\quad -2 x_{18} x_{208} +x_{19}^{2} -2 x_{19} x_{212} +x_{106}^{2} -2 x_{106} x_{216} -x_{149} x_{234} \\
&\quad +x_{153}^{2} +x_{154}^{2} +x_{162}^{2} +x_{163}^{2} +x_{169}^{2} +x_{170}^{2} +x_{175}^{2} \\
&\quad +x_{176}^{2} +x_{179}^{2} +x_{180}^{2} +x_{183}^{2} +x_{184}^{2} +x_{187}^{2} +x_{188}^{2} \\
&\quad +x_{191}^{2} +x_{192}^{2} +x_{195}^{2} +x_{196}^{2} +x_{199}^{2} +x_{200}^{2} +x_{203}^{2} \\
&\quad +x_{204}^{2} +x_{207}^{2} +x_{208}^{2} +x_{211}^{2} +x_{212}^{2} +x_{215}^{2} +x_{216}^{2}
\end{align*}

% Residual 175: remainder_budget.ratio0
\begin{align*}
R_{175} &= -x_{110} x_{235} +x_{135}
\end{align*}

% Residual 176: remainder_budget.ratio1
\begin{align*}
R_{176} &= -x_{115} x_{236} +x_{135}
\end{align*}

% Residual 177: remainder_budget.ratio2
\begin{align*}
R_{177} &= -x_{128} x_{237} +x_{135}
\end{align*}

% Residual 178: bootstrap.odd_square_radix
\begin{align*}
R_{178} &= -2 +x_{108} -4 x_{238} -4 x_{238}^{2}
\end{align*}

% Residual 179: bootstrap.guard0
\begin{align*}
R_{179} &= -1 +x_{5} +x_{5} x_{7} -x_{7} -x_{7}^{2} -x_{239}^{2} -x_{240}^{2} -x_{241} -x_{241}^{2}
\end{align*}

% Residual 180: bootstrap.guard1
\begin{align*}
R_{180} &= -2 -x_{106}^{2} -x_{154}^{2} -x_{163}^{2} +2 x_{170} -x_{170}^{2} -x_{176}^{2} -x_{180}^{2} \\
&\quad -x_{184}^{2} -x_{188}^{2} -x_{192}^{2} -x_{196}^{2} -x_{200}^{2} -x_{204}^{2} -x_{208}^{2} \\
&\quad -x_{212}^{2} +x_{216}^{2} -x_{242}^{2} -x_{243}^{2} -x_{244} -x_{244}^{2}
\end{align*}

% Residual 181: bootstrap.guard2
\begin{align*}
R_{181} &= -1 -4 x_{112}^{2} -4 x_{117}^{2} -4 x_{120}^{2} -4 x_{123}^{2} -4 x_{126}^{2} -4 x_{130}^{2} \\
&\quad -4 x_{133}^{2} +x_{135}^{2} -4 x_{137}^{2} -4 x_{141}^{2} -4 x_{144}^{2} -4 x_{147}^{2} \\
&\quad -x_{245}^{2} -x_{246}^{2} -x_{247} -x_{247}^{2}
\end{align*}
% END residuals248.tex

\endgroup

\clearpage
\begingroup
\setlength{\LTpre}{4pt}
\setlength{\LTpost}{6pt}
% BEGIN pivots.tex
\section{Complete monic-elimination pivot tables}
\label{app:monic-pivots}

Every entry $j\mapsto v$ selects the source residual $R_j$ and source
coordinate $x_v$ of the printed $248$-variable equation.  Both indices are
zero based: $0\le j<182$ and $0\le v<248$.  Write the selected row uniquely
as $R_j=\varepsilon_vx_v+f_v$ with $\varepsilon_v=\pm1$, and replace
$x_v$ by $-\varepsilon_vf_v$.  The substitution dependency graph is acyclic;
apply the definitions recursively, evaluating dependencies before their
users (a reverse topological order for the graph defined in the theorem). Retain, in
increasing source-index order, every coordinate not selected as a target.
The new polynomial is the sum of the squares of all unselected residuals
after these substitutions.  Thus these finite tables and the complete
source residual list specify all eight polynomials without a search rule
or unspecified compiler choice.  The machine-readable certificates and
literal polynomial DAGs supply the same definitions independently.

\subsection{225 integer variables, degree 6}
There are $23$ selected coordinates and $159$ remaining residuals.
\begingroup
\small
\setlength{\tabcolsep}{4pt}
\begin{longtable}{*{6}{r@{\,$\mapsto$\,}r}}
\hline
\multicolumn{12}{c}{Source residual $j$ $\mapsto$ eliminated coordinate $v$}\\
\hline
\endfirsthead
\hline
\multicolumn{12}{c}{Source residual $j$ $\mapsto$ eliminated coordinate $v$ (continued)}\\
\hline
\endhead
2 & 10 & 56 & 72 & 69 & 85 & 73 & 88 & 74 & 90 & 77 & 93 \\
79 & 95 & 81 & 97 & 83 & 99 & 86 & 102 & 87 & 103 & 91 & 106 \\
92 & 113 & 100 & 138 & 110 & 156 & 112 & 158 & 114 & 160 & 118 & 165 \\
120 & 167 & 125 & 173 & 160 & 148 & 168 & 233 & 170 & 225 & \multicolumn{2}{c}{} \\
\hline
\end{longtable}
\endgroup

\subsection{153 integer variables, degree 8}
There are $95$ selected coordinates and $87$ remaining residuals.
\begingroup
\small
\setlength{\tabcolsep}{4pt}
\begin{longtable}{*{6}{r@{\,$\mapsto$\,}r}}
\hline
\multicolumn{12}{c}{Source residual $j$ $\mapsto$ eliminated coordinate $v$}\\
\hline
\endfirsthead
\hline
\multicolumn{12}{c}{Source residual $j$ $\mapsto$ eliminated coordinate $v$ (continued)}\\
\hline
\endhead
1 & 6 & 2 & 10 & 3 & 1 & 5 & 19 & 8 & 24 & 11 & 27 \\
13 & 29 & 15 & 31 & 17 & 33 & 18 & 34 & 20 & 36 & 23 & 39 \\
25 & 41 & 26 & 42 & 30 & 46 & 31 & 47 & 34 & 50 & 37 & 53 \\
40 & 56 & 42 & 58 & 45 & 61 & 47 & 63 & 48 & 64 & 51 & 67 \\
54 & 70 & 56 & 72 & 58 & 74 & 60 & 76 & 62 & 78 & 64 & 80 \\
66 & 82 & 68 & 84 & 69 & 85 & 72 & 87 & 73 & 88 & 74 & 90 \\
77 & 93 & 79 & 95 & 81 & 97 & 83 & 99 & 84 & 100 & 86 & 102 \\
88 & 104 & 90 & 105 & 91 & 106 & 92 & 113 & 94 & 117 & 95 & 120 \\
96 & 123 & 97 & 126 & 98 & 130 & 99 & 133 & 101 & 137 & 102 & 141 \\
103 & 144 & 104 & 147 & 106 & 152 & 108 & 155 & 110 & 156 & 112 & 158 \\
114 & 160 & 116 & 164 & 119 & 162 & 121 & 153 & 122 & 171 & 124 & 169 \\
125 & 173 & 127 & 177 & 129 & 175 & 130 & 181 & 132 & 179 & 133 & 185 \\
135 & 183 & 136 & 189 & 138 & 187 & 139 & 193 & 141 & 191 & 142 & 197 \\
144 & 195 & 145 & 201 & 147 & 199 & 148 & 205 & 150 & 203 & 151 & 209 \\
153 & 207 & 154 & 213 & 156 & 211 & 157 & 217 & 162 & 0 & 164 & 231 \\
168 & 233 & 170 & 225 & 172 & 215 & 173 & 224 & 175 & 135 & \multicolumn{2}{c}{} \\
\hline
\end{longtable}
\endgroup

\subsection{144 integer variables, degree 12}
There are $104$ selected coordinates and $78$ remaining residuals.
\begingroup
\small
\setlength{\tabcolsep}{4pt}
\begin{longtable}{*{6}{r@{\,$\mapsto$\,}r}}
\hline
\multicolumn{12}{c}{Source residual $j$ $\mapsto$ eliminated coordinate $v$}\\
\hline
\endfirsthead
\hline
\multicolumn{12}{c}{Source residual $j$ $\mapsto$ eliminated coordinate $v$ (continued)}\\
\hline
\endhead
0 & 4 & 1 & 6 & 2 & 10 & 3 & 1 & 5 & 19 & 8 & 24 \\
13 & 29 & 16 & 32 & 17 & 33 & 18 & 34 & 20 & 36 & 23 & 39 \\
26 & 42 & 28 & 44 & 30 & 46 & 31 & 47 & 34 & 50 & 37 & 53 \\
40 & 56 & 42 & 58 & 45 & 61 & 47 & 63 & 48 & 64 & 51 & 67 \\
56 & 72 & 57 & 73 & 58 & 74 & 59 & 75 & 60 & 76 & 61 & 77 \\
62 & 78 & 63 & 79 & 64 & 80 & 66 & 82 & 67 & 83 & 68 & 84 \\
69 & 85 & 71 & 86 & 72 & 87 & 73 & 88 & 74 & 90 & 77 & 93 \\
79 & 95 & 81 & 97 & 83 & 99 & 84 & 100 & 86 & 102 & 88 & 104 \\
90 & 105 & 91 & 106 & 92 & 113 & 93 & 112 & 94 & 117 & 95 & 120 \\
96 & 123 & 97 & 126 & 98 & 130 & 99 & 133 & 100 & 138 & 101 & 137 \\
102 & 141 & 103 & 144 & 104 & 147 & 105 & 151 & 106 & 152 & 108 & 155 \\
110 & 156 & 112 & 158 & 114 & 160 & 116 & 164 & 119 & 162 & 121 & 153 \\
122 & 171 & 124 & 114 & 125 & 173 & 126 & 115 & 127 & 177 & 129 & 175 \\
130 & 181 & 132 & 179 & 133 & 185 & 135 & 183 & 136 & 189 & 138 & 187 \\
139 & 193 & 141 & 191 & 142 & 197 & 144 & 195 & 145 & 201 & 147 & 199 \\
148 & 205 & 150 & 203 & 151 & 209 & 153 & 207 & 154 & 213 & 156 & 211 \\
157 & 217 & 159 & 145 & 162 & 0 & 164 & 231 & 168 & 233 & 170 & 225 \\
172 & 215 & 173 & 224 & \multicolumn{2}{c}{} & \multicolumn{2}{c}{} & \multicolumn{2}{c}{} & \multicolumn{2}{c}{} \\
\hline
\end{longtable}
\endgroup

\subsection{125 integer variables, degree 16}
There are $123$ selected coordinates and $59$ remaining residuals.
\begingroup
\small
\setlength{\tabcolsep}{4pt}
\begin{longtable}{*{6}{r@{\,$\mapsto$\,}r}}
\hline
\multicolumn{12}{c}{Source residual $j$ $\mapsto$ eliminated coordinate $v$}\\
\hline
\endfirsthead
\hline
\multicolumn{12}{c}{Source residual $j$ $\mapsto$ eliminated coordinate $v$ (continued)}\\
\hline
\endhead
0 & 4 & 1 & 6 & 2 & 10 & 3 & 1 & 5 & 19 & 7 & 23 \\
8 & 24 & 13 & 29 & 16 & 32 & 17 & 33 & 18 & 34 & 20 & 36 \\
22 & 38 & 23 & 39 & 25 & 41 & 26 & 42 & 28 & 44 & 29 & 45 \\
30 & 46 & 31 & 47 & 33 & 49 & 34 & 50 & 35 & 51 & 37 & 53 \\
38 & 54 & 40 & 56 & 42 & 58 & 43 & 59 & 45 & 61 & 47 & 63 \\
48 & 64 & 49 & 65 & 51 & 67 & 52 & 68 & 54 & 70 & 56 & 72 \\
57 & 73 & 58 & 74 & 59 & 75 & 60 & 76 & 61 & 77 & 62 & 78 \\
63 & 79 & 64 & 80 & 66 & 82 & 67 & 83 & 68 & 84 & 69 & 85 \\
71 & 86 & 72 & 87 & 73 & 88 & 74 & 90 & 75 & 91 & 76 & 92 \\
77 & 93 & 79 & 95 & 81 & 97 & 83 & 99 & 84 & 100 & 86 & 102 \\
87 & 103 & 88 & 104 & 90 & 105 & 91 & 106 & 92 & 113 & 93 & 112 \\
94 & 117 & 95 & 120 & 96 & 123 & 97 & 126 & 98 & 130 & 99 & 133 \\
100 & 138 & 101 & 137 & 102 & 141 & 103 & 144 & 104 & 147 & 105 & 151 \\
106 & 152 & 108 & 155 & 110 & 156 & 112 & 158 & 114 & 160 & 116 & 164 \\
118 & 165 & 119 & 162 & 120 & 167 & 121 & 153 & 122 & 171 & 124 & 114 \\
125 & 173 & 126 & 115 & 127 & 177 & 129 & 175 & 130 & 181 & 132 & 179 \\
133 & 185 & 135 & 183 & 136 & 189 & 138 & 187 & 139 & 193 & 141 & 191 \\
142 & 197 & 144 & 195 & 145 & 201 & 147 & 199 & 148 & 205 & 150 & 203 \\
151 & 209 & 153 & 207 & 154 & 213 & 156 & 211 & 157 & 217 & 159 & 145 \\
161 & 219 & 162 & 0 & 164 & 231 & 166 & 232 & 168 & 233 & 170 & 225 \\
172 & 215 & 173 & 224 & 175 & 135 & \multicolumn{2}{c}{} & \multicolumn{2}{c}{} & \multicolumn{2}{c}{} \\
\hline
\end{longtable}
\endgroup

\subsection{117 integer variables, degree 32}
There are $131$ selected coordinates and $51$ remaining residuals.
\begingroup
\small
\setlength{\tabcolsep}{4pt}
\begin{longtable}{*{6}{r@{\,$\mapsto$\,}r}}
\hline
\multicolumn{12}{c}{Source residual $j$ $\mapsto$ eliminated coordinate $v$}\\
\hline
\endfirsthead
\hline
\multicolumn{12}{c}{Source residual $j$ $\mapsto$ eliminated coordinate $v$ (continued)}\\
\hline
\endhead
0 & 4 & 1 & 6 & 2 & 10 & 3 & 1 & 5 & 19 & 6 & 22 \\
7 & 23 & 8 & 24 & 13 & 29 & 15 & 31 & 18 & 34 & 20 & 36 \\
22 & 38 & 23 & 39 & 26 & 42 & 28 & 44 & 29 & 45 & 30 & 46 \\
31 & 47 & 33 & 49 & 34 & 50 & 35 & 51 & 37 & 53 & 38 & 54 \\
40 & 56 & 42 & 58 & 43 & 59 & 45 & 61 & 47 & 63 & 48 & 64 \\
49 & 65 & 50 & 66 & 51 & 67 & 52 & 68 & 54 & 70 & 56 & 72 \\
57 & 73 & 58 & 74 & 59 & 75 & 60 & 76 & 61 & 77 & 62 & 78 \\
63 & 79 & 64 & 80 & 65 & 81 & 66 & 82 & 67 & 83 & 68 & 84 \\
69 & 85 & 70 & 89 & 71 & 86 & 72 & 87 & 73 & 88 & 74 & 90 \\
75 & 91 & 76 & 92 & 77 & 93 & 78 & 94 & 79 & 95 & 80 & 96 \\
81 & 97 & 83 & 99 & 84 & 100 & 86 & 102 & 87 & 103 & 88 & 104 \\
90 & 105 & 91 & 106 & 92 & 113 & 93 & 112 & 94 & 117 & 95 & 120 \\
96 & 123 & 97 & 126 & 98 & 130 & 99 & 133 & 100 & 138 & 101 & 137 \\
102 & 141 & 103 & 144 & 104 & 147 & 105 & 151 & 106 & 152 & 108 & 155 \\
110 & 156 & 112 & 158 & 113 & 20 & 114 & 160 & 116 & 164 & 118 & 165 \\
119 & 162 & 120 & 167 & 121 & 153 & 122 & 171 & 124 & 114 & 125 & 173 \\
126 & 115 & 127 & 177 & 129 & 175 & 130 & 181 & 132 & 179 & 133 & 185 \\
135 & 183 & 136 & 189 & 138 & 187 & 139 & 193 & 141 & 191 & 142 & 197 \\
144 & 195 & 145 & 201 & 147 & 199 & 148 & 205 & 150 & 203 & 151 & 209 \\
153 & 207 & 154 & 213 & 156 & 211 & 157 & 217 & 159 & 145 & 161 & 219 \\
162 & 0 & 163 & 2 & 164 & 231 & 166 & 232 & 168 & 233 & 169 & 222 \\
170 & 225 & 172 & 215 & 173 & 224 & 175 & 135 & 178 & 108 & \multicolumn{2}{c}{} \\
\hline
\end{longtable}
\endgroup

\subsection{111 integer variables, degree 64}
There are $137$ selected coordinates and $45$ remaining residuals.
\begingroup
\small
\setlength{\tabcolsep}{4pt}
\begin{longtable}{*{6}{r@{\,$\mapsto$\,}r}}
\hline
\multicolumn{12}{c}{Source residual $j$ $\mapsto$ eliminated coordinate $v$}\\
\hline
\endfirsthead
\hline
\multicolumn{12}{c}{Source residual $j$ $\mapsto$ eliminated coordinate $v$ (continued)}\\
\hline
\endhead
0 & 4 & 1 & 6 & 2 & 10 & 3 & 1 & 5 & 19 & 7 & 23 \\
8 & 24 & 11 & 27 & 13 & 29 & 16 & 32 & 17 & 33 & 18 & 34 \\
19 & 35 & 20 & 36 & 22 & 38 & 23 & 39 & 25 & 41 & 26 & 42 \\
28 & 44 & 29 & 45 & 30 & 46 & 31 & 47 & 33 & 49 & 34 & 50 \\
35 & 51 & 37 & 53 & 38 & 54 & 40 & 56 & 42 & 58 & 43 & 59 \\
44 & 60 & 45 & 61 & 47 & 63 & 48 & 64 & 49 & 65 & 50 & 66 \\
51 & 67 & 52 & 68 & 54 & 70 & 56 & 72 & 57 & 73 & 58 & 74 \\
59 & 75 & 60 & 76 & 61 & 77 & 62 & 78 & 63 & 79 & 64 & 80 \\
65 & 81 & 66 & 82 & 67 & 83 & 68 & 84 & 69 & 85 & 70 & 89 \\
71 & 86 & 72 & 87 & 73 & 88 & 74 & 90 & 75 & 91 & 76 & 92 \\
77 & 93 & 78 & 94 & 79 & 95 & 80 & 96 & 81 & 97 & 82 & 98 \\
83 & 99 & 84 & 100 & 86 & 102 & 87 & 103 & 88 & 104 & 90 & 105 \\
91 & 106 & 92 & 113 & 93 & 112 & 94 & 117 & 95 & 120 & 96 & 123 \\
97 & 126 & 98 & 130 & 99 & 133 & 100 & 138 & 101 & 137 & 102 & 141 \\
103 & 144 & 104 & 147 & 105 & 151 & 106 & 152 & 108 & 155 & 110 & 156 \\
112 & 158 & 113 & 20 & 114 & 160 & 116 & 164 & 118 & 165 & 119 & 162 \\
120 & 167 & 121 & 153 & 122 & 171 & 124 & 114 & 125 & 173 & 126 & 115 \\
127 & 177 & 129 & 175 & 130 & 181 & 132 & 179 & 133 & 185 & 135 & 183 \\
136 & 189 & 138 & 187 & 139 & 193 & 141 & 191 & 142 & 197 & 144 & 195 \\
145 & 201 & 147 & 199 & 148 & 205 & 150 & 203 & 151 & 209 & 153 & 207 \\
154 & 213 & 156 & 211 & 157 & 217 & 159 & 145 & 161 & 219 & 162 & 0 \\
163 & 2 & 164 & 231 & 166 & 232 & 168 & 233 & 169 & 222 & 170 & 225 \\
171 & 149 & 172 & 215 & 173 & 224 & 175 & 135 & 178 & 108 & \multicolumn{2}{c}{} \\
\hline
\end{longtable}
\endgroup

\subsection{102 integer variables, degree 256}
There are $146$ selected coordinates and $36$ remaining residuals.
\begingroup
\small
\setlength{\tabcolsep}{4pt}
\begin{longtable}{*{6}{r@{\,$\mapsto$\,}r}}
\hline
\multicolumn{12}{c}{Source residual $j$ $\mapsto$ eliminated coordinate $v$}\\
\hline
\endfirsthead
\hline
\multicolumn{12}{c}{Source residual $j$ $\mapsto$ eliminated coordinate $v$ (continued)}\\
\hline
\endhead
0 & 4 & 1 & 6 & 2 & 10 & 3 & 1 & 5 & 19 & 7 & 23 \\
8 & 24 & 11 & 27 & 13 & 29 & 14 & 30 & 15 & 31 & 16 & 32 \\
17 & 33 & 18 & 34 & 19 & 35 & 20 & 36 & 21 & 37 & 22 & 38 \\
23 & 39 & 25 & 41 & 26 & 42 & 27 & 43 & 28 & 44 & 29 & 45 \\
30 & 46 & 31 & 47 & 32 & 48 & 33 & 49 & 34 & 50 & 35 & 51 \\
36 & 52 & 37 & 53 & 38 & 54 & 40 & 56 & 42 & 58 & 43 & 59 \\
44 & 60 & 45 & 61 & 47 & 63 & 48 & 64 & 49 & 65 & 50 & 66 \\
51 & 67 & 52 & 68 & 54 & 70 & 55 & 71 & 56 & 72 & 57 & 73 \\
58 & 74 & 59 & 75 & 60 & 76 & 61 & 77 & 62 & 78 & 63 & 79 \\
64 & 80 & 65 & 81 & 66 & 82 & 67 & 83 & 68 & 84 & 69 & 85 \\
70 & 89 & 71 & 86 & 72 & 87 & 73 & 88 & 74 & 90 & 75 & 91 \\
76 & 92 & 77 & 93 & 78 & 94 & 79 & 95 & 80 & 96 & 81 & 97 \\
82 & 98 & 83 & 99 & 84 & 100 & 85 & 101 & 86 & 102 & 87 & 103 \\
88 & 104 & 90 & 105 & 91 & 106 & 92 & 113 & 93 & 112 & 94 & 117 \\
95 & 120 & 96 & 123 & 97 & 126 & 98 & 130 & 99 & 133 & 100 & 138 \\
101 & 137 & 102 & 141 & 103 & 144 & 104 & 147 & 105 & 151 & 106 & 152 \\
108 & 155 & 110 & 156 & 112 & 158 & 113 & 20 & 114 & 160 & 116 & 164 \\
118 & 165 & 119 & 162 & 120 & 167 & 121 & 153 & 122 & 171 & 124 & 114 \\
125 & 173 & 126 & 115 & 127 & 177 & 129 & 175 & 130 & 181 & 132 & 179 \\
133 & 185 & 135 & 183 & 136 & 189 & 138 & 187 & 139 & 193 & 141 & 191 \\
142 & 197 & 144 & 195 & 145 & 201 & 147 & 199 & 148 & 205 & 150 & 203 \\
151 & 209 & 153 & 207 & 154 & 213 & 156 & 211 & 157 & 217 & 159 & 145 \\
160 & 148 & 161 & 219 & 162 & 0 & 163 & 2 & 164 & 231 & 166 & 232 \\
168 & 233 & 169 & 222 & 170 & 225 & 171 & 149 & 172 & 215 & 173 & 224 \\
175 & 135 & 178 & 108 & \multicolumn{2}{c}{} & \multicolumn{2}{c}{} & \multicolumn{2}{c}{} & \multicolumn{2}{c}{} \\
\hline
\end{longtable}
\endgroup

\subsection{94 integer variables, degree 65,540}
There are $154$ selected coordinates and $28$ remaining residuals.
\begingroup
\small
\setlength{\tabcolsep}{4pt}
\begin{longtable}{*{6}{r@{\,$\mapsto$\,}r}}
\hline
\multicolumn{12}{c}{Source residual $j$ $\mapsto$ eliminated coordinate $v$}\\
\hline
\endfirsthead
\hline
\multicolumn{12}{c}{Source residual $j$ $\mapsto$ eliminated coordinate $v$ (continued)}\\
\hline
\endhead
0 & 4 & 1 & 6 & 2 & 10 & 3 & 1 & 5 & 19 & 6 & 22 \\
7 & 23 & 8 & 24 & 9 & 25 & 10 & 26 & 11 & 27 & 12 & 28 \\
13 & 29 & 14 & 30 & 15 & 31 & 16 & 32 & 17 & 33 & 18 & 34 \\
19 & 35 & 20 & 36 & 21 & 37 & 22 & 38 & 23 & 39 & 24 & 40 \\
25 & 41 & 26 & 42 & 27 & 43 & 28 & 44 & 29 & 45 & 30 & 46 \\
31 & 47 & 32 & 48 & 33 & 49 & 34 & 50 & 35 & 51 & 36 & 52 \\
37 & 53 & 38 & 54 & 39 & 55 & 40 & 56 & 42 & 58 & 43 & 59 \\
44 & 60 & 45 & 61 & 46 & 62 & 47 & 63 & 48 & 64 & 49 & 65 \\
50 & 66 & 51 & 67 & 52 & 68 & 53 & 69 & 54 & 70 & 55 & 71 \\
56 & 72 & 57 & 73 & 58 & 74 & 59 & 75 & 60 & 76 & 61 & 77 \\
62 & 78 & 63 & 79 & 64 & 80 & 65 & 81 & 66 & 82 & 67 & 83 \\
68 & 84 & 69 & 85 & 70 & 89 & 71 & 86 & 72 & 87 & 73 & 88 \\
74 & 90 & 75 & 91 & 76 & 92 & 77 & 93 & 78 & 94 & 79 & 95 \\
80 & 96 & 81 & 97 & 82 & 98 & 83 & 99 & 84 & 100 & 85 & 101 \\
86 & 102 & 87 & 103 & 88 & 104 & 90 & 105 & 91 & 106 & 92 & 113 \\
93 & 112 & 94 & 117 & 95 & 120 & 96 & 123 & 97 & 126 & 98 & 130 \\
99 & 133 & 100 & 138 & 101 & 137 & 102 & 141 & 103 & 144 & 104 & 147 \\
105 & 151 & 106 & 152 & 108 & 155 & 110 & 156 & 112 & 158 & 113 & 20 \\
114 & 160 & 116 & 164 & 118 & 165 & 119 & 162 & 120 & 167 & 121 & 153 \\
122 & 171 & 124 & 114 & 125 & 173 & 126 & 115 & 127 & 177 & 129 & 175 \\
130 & 181 & 132 & 179 & 133 & 185 & 135 & 183 & 136 & 189 & 138 & 187 \\
139 & 193 & 141 & 191 & 142 & 197 & 144 & 195 & 145 & 201 & 147 & 199 \\
148 & 205 & 150 & 203 & 151 & 209 & 153 & 207 & 154 & 213 & 156 & 211 \\
157 & 217 & 159 & 145 & 160 & 148 & 161 & 219 & 162 & 0 & 163 & 2 \\
164 & 231 & 166 & 232 & 168 & 233 & 169 & 222 & 170 & 225 & 171 & 149 \\
172 & 215 & 173 & 224 & 175 & 135 & 178 & 108 & \multicolumn{2}{c}{} & \multicolumn{2}{c}{} \\
\hline
\end{longtable}
\endgroup

% END pivots.tex

\endgroup

\clearpage
\begin{thebibliography}{99}
\bibitem{repository}
Vladimir Reshetnikov and ProveIt contributors.
\emph{ProveIt}, comparison snapshot
\texttt{3ecd6a6cf53d47519228328567f45fb7ee9781d9}, 25 October 2026.
\url{https://github.com/VladimirReshetnikov/ProveIt/tree/3ecd6a6cf53d47519228328567f45fb7ee9781d9}.

\bibitem{forms}
ProveIt contributors. \emph{Con(ZFC) in Diophantine form}, delivered
manuscript index and comparison notes, at the same snapshot.
\url{https://github.com/VladimirReshetnikov/ProveIt/tree/3ecd6a6cf53d47519228328567f45fb7ee9781d9/Computability/DiophantineConsistency/con-zfc-diophantine-forms}.

\bibitem{review}
ProveIt contributors. \emph{Latest proposal review}, at the same snapshot.
\url{https://github.com/VladimirReshetnikov/ProveIt/blob/3ecd6a6cf53d47519228328567f45fb7ee9781d9/Computability/DiophantineConsistency/PROPOSAL_REVIEW_LATEST.md}.

\bibitem{centered}
\emph{Centered Residues}, incoming archive, 25 October 2026.
Compact component, including the retained arithmetic proofs and indexed
constant recipe.  Original notices are preserved in this package.
\url{https://github.com/VladimirReshetnikov/ProveIt/blob/3ecd6a6cf53d47519228328567f45fb7ee9781d9/docs/incoming/ConZFC_Centered_Residues_2026-10-25.zip}.

\bibitem{denominator}
\emph{Residual Bases and Denominator Engineering}, incoming archive,
25 October 2026.  In particular \texttt{compact/PROOF.md}, Section 8,
records the rejected $Y=qs$ first-order approach.
\url{https://github.com/VladimirReshetnikov/ProveIt/blob/3ecd6a6cf53d47519228328567f45fb7ee9781d9/docs/incoming/ConZFC_Residual_Bases_Denominator_Engineering_2026-10-25.zip}.

\bibitem{coprime}
\emph{Coprime Products and Positive Centers}, incoming archive,
25 October 2026.  The partial-scale predecessor motivates delayed recovery
of divisibility in a rational carrier.
\url{https://github.com/VladimirReshetnikov/ProveIt/blob/3ecd6a6cf53d47519228328567f45fb7ee9781d9/docs/incoming/ConZFC_Coprime_Products_and_Positive_Centers_2026-10-25.zip}.

\bibitem{common}
\emph{Common Remainders and Radix 47}, incoming archive, 25 October 2026.
The direct component is retained unchanged in
\path{direct_quartics/upstream_common/}, with all $88$ copied files pinned.
\url{https://github.com/VladimirReshetnikov/ProveIt/blob/3ecd6a6cf53d47519228328567f45fb7ee9781d9/docs/incoming/ConZFC_Common_Remainders_Radix47_2026-10-25.zip}.

\bibitem{graphdigits}
\emph{Graph Digits}, delivered package 139, compact component.
The $87/323$ and $88/318$ points used in the comparison.
\url{https://github.com/VladimirReshetnikov/ProveIt/tree/3ecd6a6cf53d47519228328567f45fb7ee9781d9/Computability/DiophantineConsistency/con-zfc-diophantine-forms/139-graph-digits/compact}.

\bibitem{calibration}
ProveIt contributors. The retained interfaces
\texttt{KEndpoint.c222\_iff\_zf\_refutation} and
\texttt{ConstructibleChoice.zf\_refutation\_iff\_zfc\_refutation}.
See \path{direct_quartics/upstream_common/upstream/round17/architecture/proofs/inherited/calibration/}
and its exact upstream source pins.  The corresponding native-route
interfaces and provenance are in \path{compact/parent/proof_sources/}.
\end{thebibliography}

\end{document}
